%%%%%%%%%%%%%%%%%%%%%%%%%%%%%%%%%%%%%%%%%%%%%%%%%%%%%%%%%%%%
% 13.2 MATHEMATICAL BIOLOGY
% The Composition of Advanced Mathematics
%%%%%%%%%%%%%%%%%%%%%%%%%%%%%%%%%%%%%%%%%%%%%%%%%%%%%%%%%%%%

\section{Mathematical Biology}
\label{sec:mathematical-biology}

\prerequisite{
Mathematical modeling, calculus, differential equations, dynamical systems,
probability, statistics, linear algebra, numerical methods, and basic
scientific reasoning.
}

\learningobjectives{
By the end of this section, the reader should be able to:
\begin{itemize}
    \item explain the role of mathematics in biological modeling;
    \item construct population-growth models;
    \item distinguish exponential and logistic growth;
    \item analyze equilibria and stability in population models;
    \item understand harvesting models;
    \item construct predator-prey systems;
    \item analyze the Lotka--Volterra model;
    \item recognize competition and mutualism models;
    \item construct compartmental epidemic models;
    \item understand susceptible, exposed, infected, recovered, and removed
          compartments;
    \item derive basic SIR equations;
    \item understand the basic reproduction number \(R_0\);
    \item interpret epidemic thresholds;
    \item extend SIR models to SEIR and SEIRD systems;
    \item incorporate births, deaths, vaccination, and waning immunity;
    \item understand frequency-dependent and density-dependent transmission;
    \item recognize endemic equilibria;
    \item understand age-structured population models;
    \item recognize Leslie matrices;
    \item compute long-term population growth from matrix models;
    \item understand stage-structured populations;
    \item recognize continuous age-structured models conceptually;
    \item construct simple pharmacokinetic compartment models;
    \item understand elimination and transfer rates;
    \item recognize dose-response models;
    \item understand enzyme kinetics;
    \item derive the Michaelis--Menten form conceptually;
    \item understand biochemical reaction networks;
    \item recognize mass-action kinetics;
    \item formulate gene-regulation models;
    \item understand Hill functions;
    \item recognize switch-like biological responses;
    \item model neural activity conceptually;
    \item recognize excitable systems;
    \item understand reaction-diffusion systems;
    \item recognize diffusion-driven pattern formation conceptually;
    \item understand spatial epidemic and ecological models;
    \item recognize stochastic population models;
    \item understand birth-death processes;
    \item recognize branching processes conceptually;
    \item understand extinction probabilities conceptually;
    \item distinguish deterministic and stochastic biological models;
    \item understand parameter estimation in biological systems;
    \item recognize identifiability issues;
    \item perform sensitivity analysis;
    \item understand uncertainty in biological parameters;
    \item recognize model validation issues in biology;
    \item understand biological time scales;
    \item recognize multiscale biological modeling;
    \item understand network models in biology;
    \item recognize ecological, genetic, metabolic, and epidemiological
          networks;
    \item understand evolutionary dynamics conceptually;
    \item recognize replicator equations;
    \item understand frequency-dependent selection;
    \item recognize game-theoretic biological models;
    \item understand basic population genetics modeling;
    \item recognize Hardy--Weinberg equilibrium;
    \item understand mutation and selection conceptually;
    \item recognize mathematical models of cancer growth conceptually;
    \item understand tumor-growth models;
    \item recognize immune-response models conceptually;
    \item understand biological model limitations;
    \item connect mathematical biology with differential equations,
          probability, statistics, networks, optimization, and computation.
\end{itemize}
}

%%%%%%%%%%%%%%%%%%%%%%%%%%%%%%%%%%%%%%%%%%%%%%%%%%%%%%%%%%%%
\subsection{What Is Mathematical Biology?}
\label{subsec:what-is-mathematical-biology}
%%%%%%%%%%%%%%%%%%%%%%%%%%%%%%%%%%%%%%%%%%%%%%%%%%%%%%%%%%%%

Mathematical biology uses mathematical structures to describe, analyze, and
predict biological systems.

Biological systems may involve:

\[
\boxed{
\text{populations}
}
\]

\[
\boxed{
\text{disease transmission}
}
\]

\[
\boxed{
\text{ecological interactions}
}
\]

\[
\boxed{
\text{genetics}
}
\]

\[
\boxed{
\text{biochemical reactions}
}
\]

\[
\boxed{
\text{cellular processes}
}
\]

or

\[
\boxed{
\text{physiological systems}.
}
\]

Because biological systems are often nonlinear, uncertain, and
multiscale, mathematical biology draws from many areas of mathematics.

%%%%%%%%%%%%%%%%%%%%%%%%%%%%%%%%%%%%%%%%%%%%%%%%%%%%%%%%%%%%
\subsection{Why Biology Requires Models}
\label{subsec:why-biology-requires-models}
%%%%%%%%%%%%%%%%%%%%%%%%%%%%%%%%%%%%%%%%%%%%%%%%%%%%%%%%%%%%

Biological systems contain many interacting components.

Mathematical models help isolate mechanisms such as:

\[
\boxed{
\text{birth}
}
\]

\[
\boxed{
\text{death}
}
\]

\[
\boxed{
\text{infection}
}
\]

\[
\boxed{
\text{competition}
}
\]

\[
\boxed{
\text{migration}
}
\]

and

\[
\boxed{
\text{resource limitation}.
}
\]

A model can clarify which mechanisms are sufficient to explain an observed
pattern.

%%%%%%%%%%%%%%%%%%%%%%%%%%%%%%%%%%%%%%%%%%%%%%%%%%%%%%%%%%%%
\subsection{Population Growth}
\label{subsec:population-growth}
%%%%%%%%%%%%%%%%%%%%%%%%%%%%%%%%%%%%%%%%%%%%%%%%%%%%%%%%%%%%

Let

\[
P(t)
\]

denote population size.

If births occur at per-capita rate \(b\) and deaths occur at per-capita
rate \(d\), then

\[
\frac{dP}{dt}
=
bP-dP.
\]

Thus

\[
\boxed{
\frac{dP}{dt}
=
rP,
}
\]

where

\[
\boxed{
r=b-d.
}
\]

The parameter \(r\) is the net per-capita growth rate.

%%%%%%%%%%%%%%%%%%%%%%%%%%%%%%%%%%%%%%%%%%%%%%%%%%%%%%%%%%%%
\subsection{Exponential Population Growth}
\label{subsec:exponential-population-growth}
%%%%%%%%%%%%%%%%%%%%%%%%%%%%%%%%%%%%%%%%%%%%%%%%%%%%%%%%%%%%

The differential equation

\[
\frac{dP}{dt}=rP
\]

has solution

\[
\boxed{
P(t)=P_0e^{rt}.
}
\]

If

\[
r>0,
\]

the population grows.

If

\[
r<0,
\]

the population declines.

If

\[
r=0,
\]

the population remains constant.

%%%%%%%%%%%%%%%%%%%%%%%%%%%%%%%%%%%%%%%%%%%%%%%%%%%%%%%%%%%%
\subsection{Doubling Time}
\label{subsec:population-doubling-time}
%%%%%%%%%%%%%%%%%%%%%%%%%%%%%%%%%%%%%%%%%%%%%%%%%%%%%%%%%%%%

For exponential growth,

\[
P(t)=P_0e^{rt}.
\]

The doubling time \(T_d\) satisfies

\[
2P_0
=
P_0e^{rT_d}.
\]

Therefore,

\[
2=e^{rT_d}.
\]

Taking logarithms,

\[
\boxed{
T_d
=
\frac{\ln2}{r}.
}
\]

%%%%%%%%%%%%%%%%%%%%%%%%%%%%%%%%%%%%%%%%%%%%%%%%%%%%%%%%%%%%
\subsection{Worked Example: Population Doubling}
\label{subsec:worked-population-doubling}
%%%%%%%%%%%%%%%%%%%%%%%%%%%%%%%%%%%%%%%%%%%%%%%%%%%%%%%%%%%%

\begin{workedexamplebox}{Doubling Time}

Suppose a population grows at continuous rate

\[
r=0.05.
\]

Then

\[
T_d
=
\frac{\ln2}{0.05}.
\]

Thus

\[
T_d
\approx
13.86.
\]

\begin{answerbox}
\[
\boxed{
T_d\approx13.9
}
\]

time units.
\end{answerbox}

\end{workedexamplebox}

%%%%%%%%%%%%%%%%%%%%%%%%%%%%%%%%%%%%%%%%%%%%%%%%%%%%%%%%%%%%
\subsection{Resource Limitation}
\label{subsec:biological-resource-limitation}
%%%%%%%%%%%%%%%%%%%%%%%%%%%%%%%%%%%%%%%%%%%%%%%%%%%%%%%%%%%%

Exponential growth assumes unlimited resources.

Real biological populations face constraints such as:

\[
\boxed{
\text{food}
}
\]

\[
\boxed{
\text{habitat}
}
\]

\[
\boxed{
\text{space}
}
\]

and

\[
\boxed{
\text{competition}.
}
\]

A common modification is logistic growth.

%%%%%%%%%%%%%%%%%%%%%%%%%%%%%%%%%%%%%%%%%%%%%%%%%%%%%%%%%%%%
\subsection{Logistic Population Model}
\label{subsec:logistic-population-model}
%%%%%%%%%%%%%%%%%%%%%%%%%%%%%%%%%%%%%%%%%%%%%%%%%%%%%%%%%%%%

The logistic equation is

\[
\boxed{
\frac{dP}{dt}
=
rP
\left(
1-\frac{P}{K}
\right),
}
\]

where

\[
r>0
\]

is the intrinsic growth rate and

\[
K>0
\]

is the carrying capacity.

The factor

\[
1-\frac{P}{K}
\]

reduces growth as \(P\) approaches \(K\).

%%%%%%%%%%%%%%%%%%%%%%%%%%%%%%%%%%%%%%%%%%%%%%%%%%%%%%%%%%%%
\subsection{Logistic Equilibria}
\label{subsec:logistic-population-equilibria}
%%%%%%%%%%%%%%%%%%%%%%%%%%%%%%%%%%%%%%%%%%%%%%%%%%%%%%%%%%%%

Set

\[
\frac{dP}{dt}=0.
\]

Then

\[
rP
\left(
1-\frac{P}{K}
\right)
=
0.
\]

Therefore,

\[
\boxed{
P^*=0
}
\]

or

\[
\boxed{
P^*=K.
}
\]

For \(r>0\), \(P=0\) is unstable and \(P=K\) is stable for positive
populations.

%%%%%%%%%%%%%%%%%%%%%%%%%%%%%%%%%%%%%%%%%%%%%%%%%%%%%%%%%%%%
\subsection{Phase-Line Interpretation}
\label{subsec:logistic-phase-line}
%%%%%%%%%%%%%%%%%%%%%%%%%%%%%%%%%%%%%%%%%%%%%%%%%%%%%%%%%%%%

For

\[
0<P<K,
\]

we have

\[
\frac{dP}{dt}>0.
\]

For

\[
P>K,
\]

we have

\[
\frac{dP}{dt}<0.
\]

Thus trajectories move toward

\[
\boxed{
P=K.
}
\]

This gives a qualitative stability picture without solving the equation.

%%%%%%%%%%%%%%%%%%%%%%%%%%%%%%%%%%%%%%%%%%%%%%%%%%%%%%%%%%%%
\subsection{Maximum Logistic Growth}
\label{subsec:maximum-logistic-growth}
%%%%%%%%%%%%%%%%%%%%%%%%%%%%%%%%%%%%%%%%%%%%%%%%%%%%%%%%%%%%

The growth term is

\[
G(P)
=
rP
\left(
1-\frac{P}{K}
\right).
\]

Differentiate:

\[
G'(P)
=
r
\left(
1-\frac{2P}{K}
\right).
\]

Setting

\[
G'(P)=0
\]

gives

\[
\boxed{
P=\frac{K}{2}.
}
\]

At this population,

\[
\boxed{
G_{\max}
=
\frac{rK}{4}.
}
\]

%%%%%%%%%%%%%%%%%%%%%%%%%%%%%%%%%%%%%%%%%%%%%%%%%%%%%%%%%%%%
\subsection{Harvesting}
\label{subsec:population-harvesting}
%%%%%%%%%%%%%%%%%%%%%%%%%%%%%%%%%%%%%%%%%%%%%%%%%%%%%%%%%%%%

Suppose a population is harvested at constant rate \(H\).

Then

\[
\boxed{
\frac{dP}{dt}
=
rP
\left(
1-\frac{P}{K}
\right)
-
H.
}
\]

Equilibria satisfy

\[
rP
\left(
1-\frac{P}{K}
\right)
=
H.
\]

The harvest cannot sustainably exceed the maximum natural growth rate.

%%%%%%%%%%%%%%%%%%%%%%%%%%%%%%%%%%%%%%%%%%%%%%%%%%%%%%%%%%%%
\subsection{Maximum Sustainable Constant Harvest}
\label{subsec:maximum-sustainable-harvest}
%%%%%%%%%%%%%%%%%%%%%%%%%%%%%%%%%%%%%%%%%%%%%%%%%%%%%%%%%%%%

Because logistic growth is maximized at

\[
P=\frac{K}{2},
\]

the maximum growth rate is

\[
\frac{rK}{4}.
\]

Therefore a necessary threshold for a constant sustainable harvest is

\[
\boxed{
H\leq
\frac{rK}{4}.
}
\]

At the critical value, the two positive equilibria merge.

%%%%%%%%%%%%%%%%%%%%%%%%%%%%%%%%%%%%%%%%%%%%%%%%%%%%%%%%%%%%
\subsection{Predator-Prey Models}
\label{subsec:predator-prey-models}
%%%%%%%%%%%%%%%%%%%%%%%%%%%%%%%%%%%%%%%%%%%%%%%%%%%%%%%%%%%%

Let

\[
x(t)
=
\text{prey population},
\]

and

\[
y(t)
=
\text{predator population}.
\]

A classical model is the Lotka--Volterra system:

\[
\boxed{
\frac{dx}{dt}
=
\alpha x-\beta xy,
}
\]

\[
\boxed{
\frac{dy}{dt}
=
\delta xy-\gamma y.
}
\]

%%%%%%%%%%%%%%%%%%%%%%%%%%%%%%%%%%%%%%%%%%%%%%%%%%%%%%%%%%%%
\subsection{Interpreting Lotka--Volterra Parameters}
\label{subsec:lotka-volterra-parameters}
%%%%%%%%%%%%%%%%%%%%%%%%%%%%%%%%%%%%%%%%%%%%%%%%%%%%%%%%%%%%

The parameters have interpretations:

\[
\alpha
=
\text{prey growth rate without predators},
\]

\[
\beta
=
\text{predation interaction coefficient},
\]

\[
\gamma
=
\text{predator death rate without prey},
\]

and

\[
\delta
=
\text{predator growth from prey consumption}.
\]

The nonlinear term

\[
xy
\]

represents predator-prey encounters.

%%%%%%%%%%%%%%%%%%%%%%%%%%%%%%%%%%%%%%%%%%%%%%%%%%%%%%%%%%%%
\subsection{Predator-Prey Equilibria}
\label{subsec:predator-prey-equilibria}
%%%%%%%%%%%%%%%%%%%%%%%%%%%%%%%%%%%%%%%%%%%%%%%%%%%%%%%%%%%%

Equilibria satisfy

\[
x(\alpha-\beta y)=0
\]

and

\[
y(\delta x-\gamma)=0.
\]

One equilibrium is

\[
\boxed{
(0,0).
}
\]

The positive equilibrium is

\[
\boxed{
\left(
\frac{\gamma}{\delta},
\frac{\alpha}{\beta}
\right).
}
\]

%%%%%%%%%%%%%%%%%%%%%%%%%%%%%%%%%%%%%%%%%%%%%%%%%%%%%%%%%%%%
\subsection{Worked Example: Predator-Prey Equilibrium}
\label{subsec:worked-predator-prey-equilibrium}
%%%%%%%%%%%%%%%%%%%%%%%%%%%%%%%%%%%%%%%%%%%%%%%%%%%%%%%%%%%%

\begin{workedexamplebox}{Positive Equilibrium}

Consider

\[
\frac{dx}{dt}
=
2x-0.5xy,
\]

\[
\frac{dy}{dt}
=
0.25xy-y.
\]

Here,

\[
\alpha=2,
\qquad
\beta=0.5,
\qquad
\delta=0.25,
\qquad
\gamma=1.
\]

Therefore,

\[
x^*
=
\frac{\gamma}{\delta}
=
\frac{1}{0.25}
=
4,
\]

and

\[
y^*
=
\frac{\alpha}{\beta}
=
\frac{2}{0.5}
=
4.
\]

\begin{answerbox}
\[
\boxed{
(x^*,y^*)=(4,4).
}
\]
\end{answerbox}

\end{workedexamplebox}

%%%%%%%%%%%%%%%%%%%%%%%%%%%%%%%%%%%%%%%%%%%%%%%%%%%%%%%%%%%%
\subsection{Predator-Prey Oscillation}
\label{subsec:predator-prey-oscillation}
%%%%%%%%%%%%%%%%%%%%%%%%%%%%%%%%%%%%%%%%%%%%%%%%%%%%%%%%%%%%

The classical Lotka--Volterra model can produce cyclic trajectories.

A typical interpretation is:

\[
\boxed{
\text{prey increase}
}
\]

\[
\longrightarrow
\]

\[
\boxed{
\text{predators increase}
}
\]

\[
\longrightarrow
\]

\[
\boxed{
\text{prey decrease}
}
\]

\[
\longrightarrow
\]

\[
\boxed{
\text{predators decrease}
}
\]

\[
\longrightarrow
\]

\[
\boxed{
\text{prey recover}.
}
\]

The idealized classical model produces neutrally stable closed orbits rather
than damped attraction to the coexistence equilibrium.

%%%%%%%%%%%%%%%%%%%%%%%%%%%%%%%%%%%%%%%%%%%%%%%%%%%%%%%%%%%%
\subsection{Competition Models}
\label{subsec:biological-competition-models}
%%%%%%%%%%%%%%%%%%%%%%%%%%%%%%%%%%%%%%%%%%%%%%%%%%%%%%%%%%%%

Two species competing for limited resources may be modeled by

\[
\boxed{
\frac{dx}{dt}
=
r_1x
\left(
1-\frac{x+\alpha_{12}y}{K_1}
\right),
}
\]

\[
\boxed{
\frac{dy}{dt}
=
r_2y
\left(
1-\frac{y+\alpha_{21}x}{K_2}
\right).
}
\]

The coefficients

\[
\alpha_{12},
\qquad
\alpha_{21}
\]

measure interspecific competition.

%%%%%%%%%%%%%%%%%%%%%%%%%%%%%%%%%%%%%%%%%%%%%%%%%%%%%%%%%%%%
\subsection{Possible Competition Outcomes}
\label{subsec:competition-outcomes}
%%%%%%%%%%%%%%%%%%%%%%%%%%%%%%%%%%%%%%%%%%%%%%%%%%%%%%%%%%%%

Depending on parameters, competition models may produce:

\[
\boxed{
\text{stable coexistence},
}
\]

\[
\boxed{
\text{species 1 exclusion},
}
\]

\[
\boxed{
\text{species 2 exclusion},
}
\]

or

\[
\boxed{
\text{outcome depending on initial conditions}.
}
\]

Thus nonlinear interactions can determine ecological community structure.

%%%%%%%%%%%%%%%%%%%%%%%%%%%%%%%%%%%%%%%%%%%%%%%%%%%%%%%%%%%%
\subsection{Mutualism}
\label{subsec:mutualism-models}
%%%%%%%%%%%%%%%%%%%%%%%%%%%%%%%%%%%%%%%%%%%%%%%%%%%%%%%%%%%%

In mutualistic interactions, each species benefits from the other.

A conceptual system may contain positive interaction terms:

\[
\boxed{
\frac{dx}{dt}
=
f(x)+\alpha xy,
}
\]

\[
\boxed{
\frac{dy}{dt}
=
g(y)+\beta xy,
}
\]

with

\[
\alpha,\beta>0.
\]

Realistic mutualistic models usually include saturation or resource limits
to prevent unrealistic unlimited growth.

%%%%%%%%%%%%%%%%%%%%%%%%%%%%%%%%%%%%%%%%%%%%%%%%%%%%%%%%%%%%
\subsection{Epidemic Modeling}
\label{subsec:epidemic-modeling}
%%%%%%%%%%%%%%%%%%%%%%%%%%%%%%%%%%%%%%%%%%%%%%%%%%%%%%%%%%%%

Compartmental epidemic models divide a population according to disease
status.

A basic SIR model uses:

\[
\boxed{
S(t)
=
\text{susceptible},
}
\]

\[
\boxed{
I(t)
=
\text{infectious},
}
\]

and

\[
\boxed{
R(t)
=
\text{recovered or removed}.
}
\]

%%%%%%%%%%%%%%%%%%%%%%%%%%%%%%%%%%%%%%%%%%%%%%%%%%%%%%%%%%%%
\subsection{The SIR Model}
\label{subsec:sir-model}
%%%%%%%%%%%%%%%%%%%%%%%%%%%%%%%%%%%%%%%%%%%%%%%%%%%%%%%%%%%%

Assume total population

\[
N=S+I+R
\]

is constant.

A frequency-dependent SIR model is

\[
\boxed{
\frac{dS}{dt}
=
-\beta
\frac{SI}{N},
}
\]

\[
\boxed{
\frac{dI}{dt}
=
\beta
\frac{SI}{N}
-
\gamma I,
}
\]

\[
\boxed{
\frac{dR}{dt}
=
\gamma I.
}
\]

%%%%%%%%%%%%%%%%%%%%%%%%%%%%%%%%%%%%%%%%%%%%%%%%%%%%%%%%%%%%
\subsection{Interpreting the SIR Parameters}
\label{subsec:sir-parameters}
%%%%%%%%%%%%%%%%%%%%%%%%%%%%%%%%%%%%%%%%%%%%%%%%%%%%%%%%%%%%

The parameter

\[
\beta
\]

controls effective infectious contact.

The parameter

\[
\gamma
\]

is the recovery or removal rate.

The average infectious duration in the simplest exponentially distributed
duration model is

\[
\boxed{
\frac{1}{\gamma}.
}
\]

%%%%%%%%%%%%%%%%%%%%%%%%%%%%%%%%%%%%%%%%%%%%%%%%%%%%%%%%%%%%
\subsection{Population Conservation in SIR}
\label{subsec:sir-conservation}
%%%%%%%%%%%%%%%%%%%%%%%%%%%%%%%%%%%%%%%%%%%%%%%%%%%%%%%%%%%%

Add the three equations:

\[
\frac{dS}{dt}
+
\frac{dI}{dt}
+
\frac{dR}{dt}
=
0.
\]

Therefore,

\[
\boxed{
S(t)+I(t)+R(t)=N
}
\]

for all \(t\), provided the initial total is \(N\).

This is a conservation law.

%%%%%%%%%%%%%%%%%%%%%%%%%%%%%%%%%%%%%%%%%%%%%%%%%%%%%%%%%%%%
\subsection{Initial Epidemic Growth}
\label{subsec:initial-epidemic-growth}
%%%%%%%%%%%%%%%%%%%%%%%%%%%%%%%%%%%%%%%%%%%%%%%%%%%%%%%%%%%%

The infected equation is

\[
\frac{dI}{dt}
=
I
\left(
\beta\frac{S}{N}
-
\gamma
\right).
\]

Thus infections initially increase when

\[
\boxed{
\beta\frac{S}{N}
>
\gamma.
}
\]

Initially, if nearly everyone is susceptible,

\[
S\approx N,
\]

so growth occurs when

\[
\boxed{
\beta>\gamma.
}
\]

%%%%%%%%%%%%%%%%%%%%%%%%%%%%%%%%%%%%%%%%%%%%%%%%%%%%%%%%%%%%
\subsection{Basic Reproduction Number}
\label{subsec:basic-reproduction-number}
%%%%%%%%%%%%%%%%%%%%%%%%%%%%%%%%%%%%%%%%%%%%%%%%%%%%%%%%%%%%

For the basic SIR model with a fully susceptible population,

\[
\boxed{
R_0
=
\frac{\beta}{\gamma}.
}
\]

The quantity \(R_0\) represents the expected number of secondary infections
generated by a typical infectious individual introduced into an otherwise
susceptible population under the model assumptions.

%%%%%%%%%%%%%%%%%%%%%%%%%%%%%%%%%%%%%%%%%%%%%%%%%%%%%%%%%%%%
\subsection{Epidemic Threshold}
\label{subsec:epidemic-threshold}
%%%%%%%%%%%%%%%%%%%%%%%%%%%%%%%%%%%%%%%%%%%%%%%%%%%%%%%%%%%%

If

\[
\boxed{
R_0>1,
}
\]

an introduced infection can initially grow.

If

\[
\boxed{
R_0<1,
}
\]

the infection declines from sufficiently small introduction into a largely
susceptible population.

More generally, the instantaneous growth condition is

\[
\boxed{
R_0\frac{S}{N}>1.
}
\]

%%%%%%%%%%%%%%%%%%%%%%%%%%%%%%%%%%%%%%%%%%%%%%%%%%%%%%%%%%%%
\subsection{Effective Reproduction Number}
\label{subsec:effective-reproduction-number}
%%%%%%%%%%%%%%%%%%%%%%%%%%%%%%%%%%%%%%%%%%%%%%%%%%%%%%%%%%%%

A simple effective reproduction number is

\[
\boxed{
R_{\mathrm{eff}}
=
R_0
\frac{S}{N}.
}
\]

As susceptible individuals are depleted,

\[
\frac{S}{N}
\]

decreases.

Eventually,

\[
R_{\mathrm{eff}}<1,
\]

causing the infectious population to decline in the simple model.

%%%%%%%%%%%%%%%%%%%%%%%%%%%%%%%%%%%%%%%%%%%%%%%%%%%%%%%%%%%%
\subsection{Worked Example: Epidemic Threshold}
\label{subsec:worked-epidemic-threshold}
%%%%%%%%%%%%%%%%%%%%%%%%%%%%%%%%%%%%%%%%%%%%%%%%%%%%%%%%%%%%

\begin{workedexamplebox}{Computing \(R_0\)}

Suppose

\[
\beta=0.30
\]

per day and

\[
\gamma=0.10
\]

per day.

Then

\[
R_0
=
\frac{0.30}{0.10}
=
3.
\]

\begin{answerbox}
\[
\boxed{
R_0=3.
}
\]
\end{answerbox}

In the simple SIR model, the infection can initially grow in a nearly fully
susceptible population.

\end{workedexamplebox}

%%%%%%%%%%%%%%%%%%%%%%%%%%%%%%%%%%%%%%%%%%%%%%%%%%%%%%%%%%%%
\subsection{Peak Infection Condition}
\label{subsec:sir-peak-infection}
%%%%%%%%%%%%%%%%%%%%%%%%%%%%%%%%%%%%%%%%%%%%%%%%%%%%%%%%%%%%

At an interior peak of \(I(t)\),

\[
\frac{dI}{dt}=0.
\]

For \(I>0\),

\[
\beta\frac{S}{N}
-
\gamma
=
0.
\]

Therefore,

\[
\boxed{
\frac{S}{N}
=
\frac{1}{R_0}.
}
\]

The infected population reaches a turning point when the susceptible
fraction reaches this threshold.

%%%%%%%%%%%%%%%%%%%%%%%%%%%%%%%%%%%%%%%%%%%%%%%%%%%%%%%%%%%%
\subsection{SEIR Model}
\label{subsec:seir-model}
%%%%%%%%%%%%%%%%%%%%%%%%%%%%%%%%%%%%%%%%%%%%%%%%%%%%%%%%%%%%

If infection includes a latent period, introduce an exposed compartment

\[
\boxed{
E(t).
}
\]

A standard SEIR model is

\[
\boxed{
\frac{dS}{dt}
=
-\beta\frac{SI}{N},
}
\]

\[
\boxed{
\frac{dE}{dt}
=
\beta\frac{SI}{N}
-
\sigma E,
}
\]

\[
\boxed{
\frac{dI}{dt}
=
\sigma E-\gamma I,
}
\]

\[
\boxed{
\frac{dR}{dt}
=
\gamma I.
}
\]

%%%%%%%%%%%%%%%%%%%%%%%%%%%%%%%%%%%%%%%%%%%%%%%%%%%%%%%%%%%%
\subsection{SEIR Parameters}
\label{subsec:seir-parameters}
%%%%%%%%%%%%%%%%%%%%%%%%%%%%%%%%%%%%%%%%%%%%%%%%%%%%%%%%%%%%

The parameter

\[
\sigma
\]

is the rate at which exposed individuals become infectious.

In the simplest exponential-transition interpretation,

\[
\boxed{
\frac{1}{\sigma}
}
\]

is the average exposed-period duration.

The infectious duration remains approximately

\[
\boxed{
\frac{1}{\gamma}.
}
\]

%%%%%%%%%%%%%%%%%%%%%%%%%%%%%%%%%%%%%%%%%%%%%%%%%%%%%%%%%%%%
\subsection{SEIRD Model}
\label{subsec:seird-model}
%%%%%%%%%%%%%%%%%%%%%%%%%%%%%%%%%%%%%%%%%%%%%%%%%%%%%%%%%%%%

To track disease-related deaths separately, include

\[
\boxed{
D(t).
}
\]

One possible SEIRD system is

\[
\frac{dS}{dt}
=
-\beta\frac{SI}{N},
\]

\[
\frac{dE}{dt}
=
\beta\frac{SI}{N}
-
\sigma E,
\]

\[
\frac{dI}{dt}
=
\sigma E
-
\gamma I
-
\mu I,
\]

\[
\frac{dR}{dt}
=
\gamma I,
\]

\[
\boxed{
\frac{dD}{dt}
=
\mu I.
}
\]

Here \(\mu\) is a disease-mortality rate.

%%%%%%%%%%%%%%%%%%%%%%%%%%%%%%%%%%%%%%%%%%%%%%%%%%%%%%%%%%%%
\subsection{Waning Immunity}
\label{subsec:waning-immunity}
%%%%%%%%%%%%%%%%%%%%%%%%%%%%%%%%%%%%%%%%%%%%%%%%%%%%%%%%%%%%

If immunity is temporary, recovered individuals may return to susceptibility.

A waning-immunity term is

\[
\boxed{
R
\xrightarrow{\omega}
S.
}
\]

Then the susceptible and recovered equations contain

\[
+\omega R
\]

and

\[
-\omega R,
\]

respectively.

This produces an SIRS-type model.

%%%%%%%%%%%%%%%%%%%%%%%%%%%%%%%%%%%%%%%%%%%%%%%%%%%%%%%%%%%%
\subsection{Vaccination}
\label{subsec:epidemic-vaccination}
%%%%%%%%%%%%%%%%%%%%%%%%%%%%%%%%%%%%%%%%%%%%%%%%%%%%%%%%%%%%

Vaccination can be represented by moving susceptible individuals into a
protected compartment.

A simple continuous vaccination term may be

\[
\boxed{
-\nu S
}
\]

in the susceptible equation.

If vaccination produces immunity directly,

\[
\boxed{
+\nu S
}
\]

can appear in the recovered or immune compartment.

%%%%%%%%%%%%%%%%%%%%%%%%%%%%%%%%%%%%%%%%%%%%%%%%%%%%%%%%%%%%
\subsection{Herd-Immunity Threshold in the Basic Model}
\label{subsec:herd-immunity-threshold}
%%%%%%%%%%%%%%%%%%%%%%%%%%%%%%%%%%%%%%%%%%%%%%%%%%%%%%%%%%%%

In the simple SIR setting, infection fails to grow when

\[
R_0\frac{S}{N}<1.
\]

Thus the susceptible fraction must satisfy

\[
\frac{S}{N}
<
\frac1{R_0}.
\]

Therefore the immune fraction must exceed approximately

\[
\boxed{
1-\frac1{R_0}.
}
\]

This is an idealized threshold under homogeneous-mixing assumptions.

%%%%%%%%%%%%%%%%%%%%%%%%%%%%%%%%%%%%%%%%%%%%%%%%%%%%%%%%%%%%
\subsection{Frequency-Dependent Transmission}
\label{subsec:frequency-dependent-transmission}
%%%%%%%%%%%%%%%%%%%%%%%%%%%%%%%%%%%%%%%%%%%%%%%%%%%%%%%%%%%%

Frequency-dependent incidence is commonly modeled as

\[
\boxed{
\beta\frac{SI}{N}.
}
\]

The per-susceptible infection rate depends on the fraction infectious:

\[
\boxed{
\beta\frac{I}{N}.
}
\]

This form is common when contact rates do not increase proportionally with
population density.

%%%%%%%%%%%%%%%%%%%%%%%%%%%%%%%%%%%%%%%%%%%%%%%%%%%%%%%%%%%%
\subsection{Density-Dependent Transmission}
\label{subsec:density-dependent-transmission}
%%%%%%%%%%%%%%%%%%%%%%%%%%%%%%%%%%%%%%%%%%%%%%%%%%%%%%%%%%%%

Density-dependent incidence may be modeled as

\[
\boxed{
\beta SI.
}
\]

Here encounter frequency increases with population density.

The appropriate transmission function depends on the biological context.

The two formulations should not be interchanged without adjusting the
meaning and units of \(\beta\).

%%%%%%%%%%%%%%%%%%%%%%%%%%%%%%%%%%%%%%%%%%%%%%%%%%%%%%%%%%%%
\subsection{Vital Dynamics}
\label{subsec:epidemic-vital-dynamics}
%%%%%%%%%%%%%%%%%%%%%%%%%%%%%%%%%%%%%%%%%%%%%%%%%%%%%%%%%%%%

Births and natural deaths may be included in long-term epidemic models.

For example,

\[
\frac{dS}{dt}
=
\Lambda
-
\beta\frac{SI}{N}
-
\mu S,
\]

where

\[
\Lambda
=
\text{recruitment rate}
\]

and

\[
\mu
=
\text{natural mortality rate}.
\]

Vital dynamics can permit endemic equilibria.

%%%%%%%%%%%%%%%%%%%%%%%%%%%%%%%%%%%%%%%%%%%%%%%%%%%%%%%%%%%%
\subsection{Endemic Equilibrium}
\label{subsec:endemic-equilibrium}
%%%%%%%%%%%%%%%%%%%%%%%%%%%%%%%%%%%%%%%%%%%%%%%%%%%%%%%%%%%%

An endemic equilibrium is an equilibrium with

\[
\boxed{
I^*>0.
}
\]

Unlike a disease-free equilibrium, infection persists at a constant positive
level.

Endemic equilibria frequently appear when susceptible individuals are
replenished through births or waning immunity.

%%%%%%%%%%%%%%%%%%%%%%%%%%%%%%%%%%%%%%%%%%%%%%%%%%%%%%%%%%%%
\subsection{Compartment Flow Diagrams}
\label{subsec:compartment-flow-diagrams}
%%%%%%%%%%%%%%%%%%%%%%%%%%%%%%%%%%%%%%%%%%%%%%%%%%%%%%%%%%%%

A compartment model may be summarized by

\[
\boxed{
S
\longrightarrow
E
\longrightarrow
I
\longrightarrow
R.
}
\]

For SEIRD,

\[
\boxed{
S
\longrightarrow
E
\longrightarrow
I
\longrightarrow
R
}
\]

with an additional branch

\[
\boxed{
I
\longrightarrow
D.
}
\]

Each arrow represents a modeled transition rate.

%%%%%%%%%%%%%%%%%%%%%%%%%%%%%%%%%%%%%%%%%%%%%%%%%%%%%%%%%%%%
\subsection{Age-Structured Populations}
\label{subsec:age-structured-populations}
%%%%%%%%%%%%%%%%%%%%%%%%%%%%%%%%%%%%%%%%%%%%%%%%%%%%%%%%%%%%

Population growth often depends strongly on age.

Let

\[
n_i(t)
\]

denote the number of individuals in age class \(i\).

A discrete age-structured model can be written

\[
\boxed{
\mathbf{n}_{t+1}
=
L\mathbf{n}_t,
}
\]

where \(L\) is a Leslie matrix.

%%%%%%%%%%%%%%%%%%%%%%%%%%%%%%%%%%%%%%%%%%%%%%%%%%%%%%%%%%%%
\subsection{Leslie Matrix}
\label{subsec:leslie-matrix}
%%%%%%%%%%%%%%%%%%%%%%%%%%%%%%%%%%%%%%%%%%%%%%%%%%%%%%%%%%%%

A typical Leslie matrix is

\[
\boxed{
L
=
\begin{bmatrix}
f_1 & f_2 & \cdots & f_n\\
s_1 & 0 & \cdots & 0\\
0 & s_2 & \ddots & \vdots\\
\vdots & \ddots & \ddots & 0
\end{bmatrix}.
}
\]

The \(f_i\) are fertility contributions.

The \(s_i\) are survival probabilities between age classes.

%%%%%%%%%%%%%%%%%%%%%%%%%%%%%%%%%%%%%%%%%%%%%%%%%%%%%%%%%%%%
\subsection{Worked Example: Two-Class Leslie Model}
\label{subsec:worked-leslie-model}
%%%%%%%%%%%%%%%%%%%%%%%%%%%%%%%%%%%%%%%%%%%%%%%%%%%%%%%%%%%%

\begin{workedexamplebox}{One Population Projection}

Let

\[
L
=
\begin{bmatrix}
0.4 & 1.2\\
0.8 & 0
\end{bmatrix}
\]

and initial population

\[
\mathbf{n}_0
=
\begin{bmatrix}
100\\
50
\end{bmatrix}.
\]

Then

\[
\mathbf{n}_1
=
L\mathbf{n}_0.
\]

Compute:

\[
\mathbf{n}_1
=
\begin{bmatrix}
0.4(100)+1.2(50)\\
0.8(100)
\end{bmatrix}.
\]

Thus

\[
\begin{answerbox}
\[
\boxed{
\mathbf{n}_1
=
\begin{bmatrix}
100\\
80
\end{bmatrix}.
}
\]
\end{answerbox}

\end{workedexamplebox}

%%%%%%%%%%%%%%%%%%%%%%%%%%%%%%%%%%%%%%%%%%%%%%%%%%%%%%%%%%%%
\subsection{Dominant Eigenvalue}
\label{subsec:leslie-dominant-eigenvalue}
%%%%%%%%%%%%%%%%%%%%%%%%%%%%%%%%%%%%%%%%%%%%%%%%%%%%%%%%%%%%

Under suitable positivity and irreducibility conditions, the dominant
eigenvalue of the projection matrix determines asymptotic population growth.

If the dominant eigenvalue is

\[
\lambda_1,
\]

then roughly,

\[
\boxed{
\lambda_1>1
\Rightarrow
\text{long-term growth},
}
\]

\[
\boxed{
\lambda_1<1
\Rightarrow
\text{long-term decline},
}
\]

and

\[
\boxed{
\lambda_1=1
\Rightarrow
\text{asymptotically stationary size scale}.
}
\]

%%%%%%%%%%%%%%%%%%%%%%%%%%%%%%%%%%%%%%%%%%%%%%%%%%%%%%%%%%%%
\subsection{Stable Age Distribution}
\label{subsec:stable-age-distribution}
%%%%%%%%%%%%%%%%%%%%%%%%%%%%%%%%%%%%%%%%%%%%%%%%%%%%%%%%%%%%

The eigenvector associated with the dominant eigenvalue describes the stable
age distribution under standard conditions.

After many projection steps,

\[
\mathbf{n}_t
\]

tends to align with this eigenvector up to a scalar factor.

Thus eigenvalue analysis connects linear algebra with population biology.

%%%%%%%%%%%%%%%%%%%%%%%%%%%%%%%%%%%%%%%%%%%%%%%%%%%%%%%%%%%%
\subsection{Stage-Structured Models}
\label{subsec:stage-structured-models}
%%%%%%%%%%%%%%%%%%%%%%%%%%%%%%%%%%%%%%%%%%%%%%%%%%%%%%%%%%%%

Biological populations may be better classified by stage than age.

Examples include:

\[
\boxed{
\text{seed}
\rightarrow
\text{seedling}
\rightarrow
\text{adult}
}
\]

or

\[
\boxed{
\text{egg}
\rightarrow
\text{larva}
\rightarrow
\text{adult}.
}
\]

Projection matrices can model survival, development, and reproduction among
stages.

%%%%%%%%%%%%%%%%%%%%%%%%%%%%%%%%%%%%%%%%%%%%%%%%%%%%%%%%%%%%
\subsection{Continuous Age Structure}
\label{subsec:continuous-age-structure}
%%%%%%%%%%%%%%%%%%%%%%%%%%%%%%%%%%%%%%%%%%%%%%%%%%%%%%%%%%%%

Let

\[
n(a,t)
\]

denote population density of age \(a\) at time \(t\).

A continuous age-structured model may involve a PDE of the form

\[
\boxed{
\frac{\partial n}{\partial t}
+
\frac{\partial n}{\partial a}
=
-\mu(a)n.
}
\]

Births enter through a boundary condition at

\[
a=0.
\]

%%%%%%%%%%%%%%%%%%%%%%%%%%%%%%%%%%%%%%%%%%%%%%%%%%%%%%%%%%%%
\subsection{Birth Boundary Condition}
\label{subsec:age-structured-birth-boundary}
%%%%%%%%%%%%%%%%%%%%%%%%%%%%%%%%%%%%%%%%%%%%%%%%%%%%%%%%%%%%

A common birth condition is

\[
\boxed{
n(0,t)
=
\int_0^\infty
b(a)n(a,t)\,da,
}
\]

where

\[
b(a)
\]

is an age-specific fertility rate.

Thus reproduction couples all reproductive age classes back to age zero.

%%%%%%%%%%%%%%%%%%%%%%%%%%%%%%%%%%%%%%%%%%%%%%%%%%%%%%%%%%%%
\subsection{Pharmacokinetic Models}
\label{subsec:pharmacokinetic-models}
%%%%%%%%%%%%%%%%%%%%%%%%%%%%%%%%%%%%%%%%%%%%%%%%%%%%%%%%%%%%

Pharmacokinetics studies how drugs move through the body.

A simple one-compartment model may describe drug amount \(A(t)\) by

\[
\boxed{
\frac{dA}{dt}
=
-kA.
}
\]

The solution is

\[
\boxed{
A(t)=A_0e^{-kt}.
}
\]

%%%%%%%%%%%%%%%%%%%%%%%%%%%%%%%%%%%%%%%%%%%%%%%%%%%%%%%%%%%%
\subsection{Drug Half-Life}
\label{subsec:drug-half-life}
%%%%%%%%%%%%%%%%%%%%%%%%%%%%%%%%%%%%%%%%%%%%%%%%%%%%%%%%%%%%

For

\[
A(t)=A_0e^{-kt},
\]

the half-life \(T_{1/2}\) satisfies

\[
\frac{A_0}{2}
=
A_0e^{-kT_{1/2}}.
\]

Therefore,

\[
\boxed{
T_{1/2}
=
\frac{\ln2}{k}.
}
\]

%%%%%%%%%%%%%%%%%%%%%%%%%%%%%%%%%%%%%%%%%%%%%%%%%%%%%%%%%%%%
\subsection{Worked Example: Drug Elimination}
\label{subsec:worked-drug-elimination}
%%%%%%%%%%%%%%%%%%%%%%%%%%%%%%%%%%%%%%%%%%%%%%%%%%%%%%%%%%%%

\begin{workedexamplebox}{Half-Life From Elimination Rate}

Suppose

\[
k=0.2
\]

per hour.

Then

\[
T_{1/2}
=
\frac{\ln2}{0.2}.
\]

Therefore,

\begin{answerbox}
\[
\boxed{
T_{1/2}
\approx3.47\text{ hours}.
}
\]
\end{answerbox}

\end{workedexamplebox}

%%%%%%%%%%%%%%%%%%%%%%%%%%%%%%%%%%%%%%%%%%%%%%%%%%%%%%%%%%%%
\subsection{Two-Compartment Pharmacokinetics}
\label{subsec:two-compartment-pharmacokinetics}
%%%%%%%%%%%%%%%%%%%%%%%%%%%%%%%%%%%%%%%%%%%%%%%%%%%%%%%%%%%%

Let

\[
A_1(t)
\]

and

\[
A_2(t)
\]

represent drug amounts in central and peripheral compartments.

A simple model is

\[
\boxed{
\frac{dA_1}{dt}
=
-k_{12}A_1
+
k_{21}A_2
-
k_eA_1,
}
\]

\[
\boxed{
\frac{dA_2}{dt}
=
k_{12}A_1
-
k_{21}A_2.
}
\]

Here \(k_e\) is the elimination rate from the central compartment.

%%%%%%%%%%%%%%%%%%%%%%%%%%%%%%%%%%%%%%%%%%%%%%%%%%%%%%%%%%%%
\subsection{Biochemical Reaction Networks}
\label{subsec:biochemical-reaction-networks}
%%%%%%%%%%%%%%%%%%%%%%%%%%%%%%%%%%%%%%%%%%%%%%%%%%%%%%%%%%%%

Biochemical systems consist of interacting reactions among chemical
species.

A reaction may be written

\[
\boxed{
A+B
\xrightarrow{k}
C.
}
\]

Under mass-action kinetics, the reaction rate is proportional to the product
of reactant concentrations:

\[
\boxed{
v=k[A][B].
}
\]

%%%%%%%%%%%%%%%%%%%%%%%%%%%%%%%%%%%%%%%%%%%%%%%%%%%%%%%%%%%%
\subsection{Mass-Action Kinetics}
\label{subsec:mass-action-kinetics}
%%%%%%%%%%%%%%%%%%%%%%%%%%%%%%%%%%%%%%%%%%%%%%%%%%%%%%%%%%%%

For the reaction

\[
A+B
\rightarrow
C,
\]

let concentrations be

\[
a(t),b(t),c(t).
\]

Then a simple model is

\[
\boxed{
\frac{da}{dt}
=
-kab,
}
\]

\[
\boxed{
\frac{db}{dt}
=
-kab,
}
\]

\[
\boxed{
\frac{dc}{dt}
=
kab.
}
\]

%%%%%%%%%%%%%%%%%%%%%%%%%%%%%%%%%%%%%%%%%%%%%%%%%%%%%%%%%%%%
\subsection{Enzyme Kinetics}
\label{subsec:enzyme-kinetics}
%%%%%%%%%%%%%%%%%%%%%%%%%%%%%%%%%%%%%%%%%%%%%%%%%%%%%%%%%%%%

A classical enzyme reaction is

\[
\boxed{
E+S
\rightleftharpoons
ES
\rightarrow
E+P.
}
\]

Here:

\[
E
=
\text{enzyme},
\]

\[
S
=
\text{substrate},
\]

\[
ES
=
\text{enzyme-substrate complex},
\]

and

\[
P
=
\text{product}.
\]

%%%%%%%%%%%%%%%%%%%%%%%%%%%%%%%%%%%%%%%%%%%%%%%%%%%%%%%%%%%%
\subsection{Michaelis--Menten Model}
\label{subsec:michaelis-menten}
%%%%%%%%%%%%%%%%%%%%%%%%%%%%%%%%%%%%%%%%%%%%%%%%%%%%%%%%%%%%

Under standard quasi-steady-state assumptions, the reaction velocity can be
approximated by

\[
\boxed{
v
=
\frac{V_{\max}S}
{K_M+S}.
}
\]

Here

\[
V_{\max}
\]

is the maximum reaction rate and

\[
K_M
\]

is the Michaelis constant.

%%%%%%%%%%%%%%%%%%%%%%%%%%%%%%%%%%%%%%%%%%%%%%%%%%%%%%%%%%%%
\subsection{Michaelis--Menten Limits}
\label{subsec:michaelis-menten-limits}
%%%%%%%%%%%%%%%%%%%%%%%%%%%%%%%%%%%%%%%%%%%%%%%%%%%%%%%%%%%%

If

\[
S\ll K_M,
\]

then

\[
v
\approx
\frac{V_{\max}}{K_M}S.
\]

Thus the response is approximately linear.

If

\[
S\gg K_M,
\]

then

\[
\boxed{
v\approx V_{\max}.
}
\]

The enzyme becomes saturated.

%%%%%%%%%%%%%%%%%%%%%%%%%%%%%%%%%%%%%%%%%%%%%%%%%%%%%%%%%%%%
\subsection{Worked Example: Half-Maximal Enzyme Rate}
\label{subsec:worked-michaelis-half}
%%%%%%%%%%%%%%%%%%%%%%%%%%%%%%%%%%%%%%%%%%%%%%%%%%%%%%%%%%%%

\begin{workedexamplebox}{Substrate at \(K_M\)}

Set

\[
S=K_M.
\]

Then

\[
v
=
\frac{V_{\max}K_M}
{K_M+K_M}.
\]

Therefore,

\[
\begin{answerbox}
\[
\boxed{
v
=
\frac{V_{\max}}{2}.
}
\]
\end{answerbox}

Thus \(K_M\) is the substrate concentration producing half the maximum rate
in the Michaelis--Menten model.

\end{workedexamplebox}

%%%%%%%%%%%%%%%%%%%%%%%%%%%%%%%%%%%%%%%%%%%%%%%%%%%%%%%%%%%%
\subsection{Hill Functions}
\label{subsec:hill-functions}
%%%%%%%%%%%%%%%%%%%%%%%%%%%%%%%%%%%%%%%%%%%%%%%%%%%%%%%%%%%%

Biological responses often show cooperative or switch-like behavior.

A Hill function is

\[
\boxed{
H(x)
=
\frac{x^n}
{K^n+x^n}.
}
\]

Here:

\[
K
=
\text{half-saturation scale},
\]

and

\[
n
=
\text{Hill coefficient}.
\]

Larger \(n\) produces a sharper transition.

%%%%%%%%%%%%%%%%%%%%%%%%%%%%%%%%%%%%%%%%%%%%%%%%%%%%%%%%%%%%
\subsection{Gene Regulation}
\label{subsec:gene-regulation-models}
%%%%%%%%%%%%%%%%%%%%%%%%%%%%%%%%%%%%%%%%%%%%%%%%%%%%%%%%%%%%

Let

\[
m(t)
=
\text{mRNA concentration}
\]

and

\[
p(t)
=
\text{protein concentration}.
\]

A basic model is

\[
\boxed{
\frac{dm}{dt}
=
\alpha
-
\delta_m m,
}
\]

\[
\boxed{
\frac{dp}{dt}
=
\beta m
-
\delta_p p.
}
\]

Regulatory effects can be incorporated by replacing constant production
rates with Hill-type functions.

%%%%%%%%%%%%%%%%%%%%%%%%%%%%%%%%%%%%%%%%%%%%%%%%%%%%%%%%%%%%
\subsection{Negative Gene Regulation}
\label{subsec:negative-gene-regulation}
%%%%%%%%%%%%%%%%%%%%%%%%%%%%%%%%%%%%%%%%%%%%%%%%%%%%%%%%%%%%

A repressive production law may take form

\[
\boxed{
\alpha(p)
=
\frac{\alpha_{\max}}
{1+(p/K)^n}.
}
\]

As \(p\) increases, production decreases.

This creates negative feedback.

Negative feedback can stabilize concentrations and reduce fluctuations.

%%%%%%%%%%%%%%%%%%%%%%%%%%%%%%%%%%%%%%%%%%%%%%%%%%%%%%%%%%%%
\subsection{Positive Gene Regulation}
\label{subsec:positive-gene-regulation}
%%%%%%%%%%%%%%%%%%%%%%%%%%%%%%%%%%%%%%%%%%%%%%%%%%%%%%%%%%%%

An activating law may be modeled as

\[
\boxed{
\alpha(p)
=
\alpha_{\max}
\frac{p^n}
{K^n+p^n}.
}
\]

Positive feedback can create:

\[
\boxed{
\text{switch-like behavior},
}
\]

\[
\boxed{
\text{multiple equilibria},
}
\]

or

\[
\boxed{
\text{bistability}.
}
\]

%%%%%%%%%%%%%%%%%%%%%%%%%%%%%%%%%%%%%%%%%%%%%%%%%%%%%%%%%%%%
\subsection{Neural Models}
\label{subsec:neural-models}
%%%%%%%%%%%%%%%%%%%%%%%%%%%%%%%%%%%%%%%%%%%%%%%%%%%%%%%%%%%%

Mathematical neuroscience models electrical activity of neurons.

A simple membrane equation may be written

\[
\boxed{
C\frac{dV}{dt}
=
-I_{\mathrm{ion}}(V,\mathbf{w})
+
I_{\mathrm{ext}},
}
\]

where:

\[
V
=
\text{membrane potential},
\]

\[
C
=
\text{membrane capacitance},
\]

and

\[
I_{\mathrm{ext}}
=
\text{external input current}.
\]

%%%%%%%%%%%%%%%%%%%%%%%%%%%%%%%%%%%%%%%%%%%%%%%%%%%%%%%%%%%%
\subsection{Excitable Systems}
\label{subsec:excitable-systems}
%%%%%%%%%%%%%%%%%%%%%%%%%%%%%%%%%%%%%%%%%%%%%%%%%%%%%%%%%%%%

An excitable system has a stable resting state but can generate a large
temporary response when a threshold is crossed.

This behavior appears in:

\[
\boxed{
\text{neurons},
}
\]

\[
\boxed{
\text{cardiac cells},
}
\]

and

\[
\boxed{
\text{chemical signaling systems}.
}
\]

Nonlinear dynamical systems provide the mathematical framework.

%%%%%%%%%%%%%%%%%%%%%%%%%%%%%%%%%%%%%%%%%%%%%%%%%%%%%%%%%%%%
\subsection{Reaction-Diffusion Systems}
\label{subsec:biological-reaction-diffusion}
%%%%%%%%%%%%%%%%%%%%%%%%%%%%%%%%%%%%%%%%%%%%%%%%%%%%%%%%%%%%

Biological substances may both react and diffuse.

A two-species system is

\[
\boxed{
\frac{\partial u}{\partial t}
=
D_u\nabla^2u
+
f(u,v),
}
\]

\[
\boxed{
\frac{\partial v}{\partial t}
=
D_v\nabla^2v
+
g(u,v).
}
\]

The reaction terms describe local biology.

The diffusion terms describe spatial movement.

%%%%%%%%%%%%%%%%%%%%%%%%%%%%%%%%%%%%%%%%%%%%%%%%%%%%%%%%%%%%
\subsection{Pattern Formation}
\label{subsec:biological-pattern-formation}
%%%%%%%%%%%%%%%%%%%%%%%%%%%%%%%%%%%%%%%%%%%%%%%%%%%%%%%%%%%%

Reaction-diffusion systems can produce spatial patterns.

A spatially homogeneous equilibrium may be stable without diffusion yet
become unstable after diffusion is introduced under suitable conditions.

This mechanism is associated with diffusion-driven pattern formation.

Such models have been used conceptually for biological patterning.

%%%%%%%%%%%%%%%%%%%%%%%%%%%%%%%%%%%%%%%%%%%%%%%%%%%%%%%%%%%%
\subsection{Spatial Ecology}
\label{subsec:spatial-ecology}
%%%%%%%%%%%%%%%%%%%%%%%%%%%%%%%%%%%%%%%%%%%%%%%%%%%%%%%%%%%%

A spatial population model may take form

\[
\boxed{
\frac{\partial P}{\partial t}
=
D\nabla^2P
+
rP
\left(
1-\frac{P}{K}
\right).
}
\]

This combines:

\[
\boxed{
\text{diffusion}
}
\]

and

\[
\boxed{
\text{logistic growth}.
}
\]

It is a basic reaction-diffusion population model.

%%%%%%%%%%%%%%%%%%%%%%%%%%%%%%%%%%%%%%%%%%%%%%%%%%%%%%%%%%%%
\subsection{Travel and Migration Models}
\label{subsec:biological-migration-models}
%%%%%%%%%%%%%%%%%%%%%%%%%%%%%%%%%%%%%%%%%%%%%%%%%%%%%%%%%%%%

Suppose populations occupy several regions.

Let

\[
x_i(t)
\]

be the population in region \(i\).

A network migration model can include terms

\[
\boxed{
\sum_j m_{ji}x_j
-
\sum_j m_{ij}x_i.
}
\]

The first sum represents arrivals.

The second represents departures.

%%%%%%%%%%%%%%%%%%%%%%%%%%%%%%%%%%%%%%%%%%%%%%%%%%%%%%%%%%%%
\subsection{Metapopulation Models}
\label{subsec:metapopulation-models}
%%%%%%%%%%%%%%%%%%%%%%%%%%%%%%%%%%%%%%%%%%%%%%%%%%%%%%%%%%%%

A metapopulation consists of local populations connected by migration.

The state may be

\[
\boxed{
\mathbf{x}
=
(x_1,\ldots,x_n).
}
\]

Dynamics combine local biological processes with movement among patches.

Metapopulation models are important in ecology and spatial epidemiology.

%%%%%%%%%%%%%%%%%%%%%%%%%%%%%%%%%%%%%%%%%%%%%%%%%%%%%%%%%%%%
\subsection{Stochastic Birth-Death Models}
\label{subsec:stochastic-birth-death}
%%%%%%%%%%%%%%%%%%%%%%%%%%%%%%%%%%%%%%%%%%%%%%%%%%%%%%%%%%%%

For small populations, random events may matter greatly.

Suppose a population has birth rate

\[
\lambda_n
\]

and death rate

\[
\mu_n
\]

when population size is \(n\).

Then transitions are

\[
\boxed{
n
\rightarrow
n+1
}
\]

at rate \(\lambda_n\), and

\[
\boxed{
n
\rightarrow
n-1
}
\]

at rate \(\mu_n\).

%%%%%%%%%%%%%%%%%%%%%%%%%%%%%%%%%%%%%%%%%%%%%%%%%%%%%%%%%%%%
\subsection{Linear Birth-Death Process}
\label{subsec:linear-birth-death}
%%%%%%%%%%%%%%%%%%%%%%%%%%%%%%%%%%%%%%%%%%%%%%%%%%%%%%%%%%%%

A common model uses

\[
\boxed{
\lambda_n
=
\lambda n,
}
\]

\[
\boxed{
\mu_n
=
\mu n.
}
\]

Each individual contributes independently to birth and death rates.

The expected population satisfies

\[
\boxed{
\frac{d}{dt}\mathbb{E}[N(t)]
=
(\lambda-\mu)
\mathbb{E}[N(t)].
}
\]

%%%%%%%%%%%%%%%%%%%%%%%%%%%%%%%%%%%%%%%%%%%%%%%%%%%%%%%%%%%%
\subsection{Extinction}
\label{subsec:biological-extinction}
%%%%%%%%%%%%%%%%%%%%%%%%%%%%%%%%%%%%%%%%%%%%%%%%%%%%%%%%%%%%

In a stochastic population model,

\[
\boxed{
N(t)=0
}
\]

may be an absorbing state.

Even when the deterministic mean suggests growth, a small stochastic
population may go extinct because of random fluctuations.

This is a major difference between deterministic and stochastic modeling.

%%%%%%%%%%%%%%%%%%%%%%%%%%%%%%%%%%%%%%%%%%%%%%%%%%%%%%%%%%%%
\subsection{Branching Processes}
\label{subsec:branching-processes-biology}
%%%%%%%%%%%%%%%%%%%%%%%%%%%%%%%%%%%%%%%%%%%%%%%%%%%%%%%%%%%%

In a branching process, individuals independently produce random numbers of
offspring.

Let

\[
X
\]

be the offspring number of one individual.

If

\[
\boxed{
\mathbb{E}[X]\leq1,
}
\]

extinction occurs with probability \(1\) under standard nondegenerate
Galton--Watson assumptions.

If

\[
\boxed{
\mathbb{E}[X]>1,
}
\]

the extinction probability is less than \(1\).

%%%%%%%%%%%%%%%%%%%%%%%%%%%%%%%%%%%%%%%%%%%%%%%%%%%%%%%%%%%%
\subsection{Evolutionary Dynamics}
\label{subsec:evolutionary-dynamics}
%%%%%%%%%%%%%%%%%%%%%%%%%%%%%%%%%%%%%%%%%%%%%%%%%%%%%%%%%%%%

Evolution changes trait or strategy frequencies through differential
reproductive success.

Let

\[
x_i
\]

be the frequency of strategy \(i\), with

\[
\boxed{
\sum_i x_i=1.
}
\]

Fitness may depend on the current population composition.

This creates frequency-dependent selection.

%%%%%%%%%%%%%%%%%%%%%%%%%%%%%%%%%%%%%%%%%%%%%%%%%%%%%%%%%%%%
\subsection{Replicator Equation}
\label{subsec:replicator-equation}
%%%%%%%%%%%%%%%%%%%%%%%%%%%%%%%%%%%%%%%%%%%%%%%%%%%%%%%%%%%%

The replicator equation is

\[
\boxed{
\dot{x}_i
=
x_i
\left(
f_i(\mathbf{x})
-
\bar{f}(\mathbf{x})
\right),
}
\]

where

\[
f_i(\mathbf{x})
\]

is the fitness of strategy \(i\) and

\[
\boxed{
\bar{f}
=
\sum_jx_jf_j
}
\]

is the mean population fitness.

Strategies performing above average increase in frequency.

%%%%%%%%%%%%%%%%%%%%%%%%%%%%%%%%%%%%%%%%%%%%%%%%%%%%%%%%%%%%
\subsection{Worked Example: Two-Strategy Replicator Model}
\label{subsec:worked-two-strategy-replicator}
%%%%%%%%%%%%%%%%%%%%%%%%%%%%%%%%%%%%%%%%%%%%%%%%%%%%%%%%%%%%

\begin{workedexamplebox}{Constant Fitness Difference}

Suppose strategy \(A\) has fitness

\[
f_A=2
\]

and strategy \(B\) has fitness

\[
f_B=1.
\]

Let \(x\) be the frequency of \(A\).

Then

\[
\bar f
=
2x+1(1-x)
=
1+x.
\]

Thus

\[
\dot{x}
=
x
\left(
2-(1+x)
\right).
\]

Therefore,

\begin{answerbox}
\[
\boxed{
\dot{x}
=
x(1-x).
}
\]
\end{answerbox}

Strategy \(A\) increases whenever

\[
0<x<1.
\]

\end{workedexamplebox}

%%%%%%%%%%%%%%%%%%%%%%%%%%%%%%%%%%%%%%%%%%%%%%%%%%%%%%%%%%%%
\subsection{Evolutionary Game Theory}
\label{subsec:evolutionary-game-theory}
%%%%%%%%%%%%%%%%%%%%%%%%%%%%%%%%%%%%%%%%%%%%%%%%%%%%%%%%%%%%

When fitness depends on interactions, payoff matrices can determine
reproductive success.

If

\[
A
\]

is a payoff matrix and \(\mathbf{x}\) is the strategy-frequency vector, one
may write

\[
\boxed{
f_i(\mathbf{x})
=
(A\mathbf{x})_i.
}
\]

Then

\[
\boxed{
\bar f
=
\mathbf{x}^{T}A\mathbf{x}.
}
\]

This links evolutionary biology with game theory.

%%%%%%%%%%%%%%%%%%%%%%%%%%%%%%%%%%%%%%%%%%%%%%%%%%%%%%%%%%%%
\subsection{Population Genetics}
\label{subsec:population-genetics-models}
%%%%%%%%%%%%%%%%%%%%%%%%%%%%%%%%%%%%%%%%%%%%%%%%%%%%%%%%%%%%

Population genetics studies how allele frequencies change over generations.

Suppose a gene has two alleles:

\[
\boxed{
A
}
\]

and

\[
\boxed{
a.
}
\]

Let their frequencies be

\[
p
\]

and

\[
q,
\]

with

\[
\boxed{
p+q=1.
}
\]

%%%%%%%%%%%%%%%%%%%%%%%%%%%%%%%%%%%%%%%%%%%%%%%%%%%%%%%%%%%%
\subsection{Hardy--Weinberg Equilibrium}
\label{subsec:hardy-weinberg}
%%%%%%%%%%%%%%%%%%%%%%%%%%%%%%%%%%%%%%%%%%%%%%%%%%%%%%%%%%%%

Under standard idealized assumptions, genotype frequencies after random
mating are

\[
\boxed{
AA:p^2,
}
\]

\[
\boxed{
Aa:2pq,
}
\]

\[
\boxed{
aa:q^2.
}
\]

Since

\[
p+q=1,
\]

we have

\[
\boxed{
p^2+2pq+q^2=1.
}
\]

%%%%%%%%%%%%%%%%%%%%%%%%%%%%%%%%%%%%%%%%%%%%%%%%%%%%%%%%%%%%
\subsection{Worked Example: Hardy--Weinberg Frequencies}
\label{subsec:worked-hardy-weinberg}
%%%%%%%%%%%%%%%%%%%%%%%%%%%%%%%%%%%%%%%%%%%%%%%%%%%%%%%%%%%%

\begin{workedexamplebox}{Genotype Frequencies}

Suppose

\[
p=0.7,
\qquad
q=0.3.
\]

Then

\[
p^2=0.49,
\]

\[
2pq
=
2(0.7)(0.3)
=
0.42,
\]

and

\[
q^2=0.09.
\]

\begin{answerbox}
\[
\boxed{
AA=0.49,
\qquad
Aa=0.42,
\qquad
aa=0.09.
}
\]
\end{answerbox}

\end{workedexamplebox}

%%%%%%%%%%%%%%%%%%%%%%%%%%%%%%%%%%%%%%%%%%%%%%%%%%%%%%%%%%%%
\subsection{Selection}
\label{subsec:population-genetic-selection}
%%%%%%%%%%%%%%%%%%%%%%%%%%%%%%%%%%%%%%%%%%%%%%%%%%%%%%%%%%%%

Suppose genotype fitnesses are

\[
w_{AA},
\qquad
w_{Aa},
\qquad
w_{aa}.
\]

Selection changes genotype contributions to the next generation.

Allele-frequency recursion can then be derived from weighted genotype
frequencies.

Thus evolution can be modeled as a dynamical system on frequencies.

%%%%%%%%%%%%%%%%%%%%%%%%%%%%%%%%%%%%%%%%%%%%%%%%%%%%%%%%%%%%
\subsection{Mutation}
\label{subsec:population-genetic-mutation}
%%%%%%%%%%%%%%%%%%%%%%%%%%%%%%%%%%%%%%%%%%%%%%%%%%%%%%%%%%%%

Suppose allele \(A\) mutates to \(a\) with probability \(\mu\) and \(a\)
mutates to \(A\) with probability \(\nu\).

Then a simple allele-frequency update is

\[
\boxed{
p_{t+1}
=
(1-\mu)p_t
+
\nu(1-p_t),
}
\]

before incorporating other evolutionary forces.

Mutation introduces new variation and prevents some absorbing states.

%%%%%%%%%%%%%%%%%%%%%%%%%%%%%%%%%%%%%%%%%%%%%%%%%%%%%%%%%%%%
\subsection{Cancer Growth Models}
\label{subsec:cancer-growth-models}
%%%%%%%%%%%%%%%%%%%%%%%%%%%%%%%%%%%%%%%%%%%%%%%%%%%%%%%%%%%%

Tumor growth may be modeled using:

\[
\boxed{
\text{exponential growth},
}
\]

\[
\boxed{
\text{logistic growth},
}
\]

or more specialized growth laws.

A simple logistic tumor model is

\[
\boxed{
\frac{dV}{dt}
=
rV
\left(
1-\frac{V}{K}
\right),
}
\]

where \(V(t)\) is tumor volume.

%%%%%%%%%%%%%%%%%%%%%%%%%%%%%%%%%%%%%%%%%%%%%%%%%%%%%%%%%%%%
\subsection{Gompertz Growth}
\label{subsec:gompertz-growth}
%%%%%%%%%%%%%%%%%%%%%%%%%%%%%%%%%%%%%%%%%%%%%%%%%%%%%%%%%%%%

A Gompertz growth model is

\[
\boxed{
\frac{dV}{dt}
=
rV
\ln
\left(
\frac{K}{V}
\right).
}
\]

The per-capita growth rate decreases approximately logarithmically as
population or tumor size increases.

This model is frequently used as an alternative to logistic growth.

%%%%%%%%%%%%%%%%%%%%%%%%%%%%%%%%%%%%%%%%%%%%%%%%%%%%%%%%%%%%
\subsection{Treatment Models}
\label{subsec:biological-treatment-models}
%%%%%%%%%%%%%%%%%%%%%%%%%%%%%%%%%%%%%%%%%%%%%%%%%%%%%%%%%%%%

Treatment can be incorporated as a removal or mortality term.

For example,

\[
\boxed{
\frac{dV}{dt}
=
rV
\left(
1-\frac{V}{K}
\right)
-
u(t)V,
}
\]

where

\[
u(t)
\]

represents treatment intensity.

This connects biological modeling with optimal control.

%%%%%%%%%%%%%%%%%%%%%%%%%%%%%%%%%%%%%%%%%%%%%%%%%%%%%%%%%%%%
\subsection{Immune Response Models}
\label{subsec:immune-response-models}
%%%%%%%%%%%%%%%%%%%%%%%%%%%%%%%%%%%%%%%%%%%%%%%%%%%%%%%%%%%%

Let

\[
P(t)
=
\text{pathogen population}
\]

and

\[
E(t)
=
\text{immune effector population}.
\]

A conceptual model may take form

\[
\boxed{
\frac{dP}{dt}
=
rP-kEP,
}
\]

\[
\boxed{
\frac{dE}{dt}
=
aP-bE.
}
\]

Such models represent interactions between pathogen growth and immune
suppression.

%%%%%%%%%%%%%%%%%%%%%%%%%%%%%%%%%%%%%%%%%%%%%%%%%%%%%%%%%%%%
\subsection{Biological Networks}
\label{subsec:biological-networks}
%%%%%%%%%%%%%%%%%%%%%%%%%%%%%%%%%%%%%%%%%%%%%%%%%%%%%%%%%%%%

Many biological systems have network structure.

Examples include:

\[
\boxed{
\text{gene-regulatory networks},
}
\]

\[
\boxed{
\text{metabolic networks},
}
\]

\[
\boxed{
\text{neural networks},
}
\]

\[
\boxed{
\text{food webs},
}
\]

and

\[
\boxed{
\text{contact networks}.
}
\]

Nodes represent biological entities and edges represent interactions.

%%%%%%%%%%%%%%%%%%%%%%%%%%%%%%%%%%%%%%%%%%%%%%%%%%%%%%%%%%%%
\subsection{Contact Networks and Epidemics}
\label{subsec:contact-network-epidemics}
%%%%%%%%%%%%%%%%%%%%%%%%%%%%%%%%%%%%%%%%%%%%%%%%%%%%%%%%%%%%

Homogeneous epidemic models assume that individuals mix uniformly.

A contact-network model instead represents individuals or groups as nodes.

Infection can travel only along edges.

Thus network structure can alter:

\[
\boxed{
\text{epidemic thresholds},
}
\]

\[
\boxed{
\text{outbreak sizes},
}
\]

and

\[
\boxed{
\text{intervention effectiveness}.
}
\]

%%%%%%%%%%%%%%%%%%%%%%%%%%%%%%%%%%%%%%%%%%%%%%%%%%%%%%%%%%%%
\subsection{Multiscale Biological Modeling}
\label{subsec:multiscale-biological-modeling}
%%%%%%%%%%%%%%%%%%%%%%%%%%%%%%%%%%%%%%%%%%%%%%%%%%%%%%%%%%%%

Biological phenomena span many scales.

For example:

\[
\boxed{
\text{genes}
}
\]

\[
\longrightarrow
\]

\[
\boxed{
\text{proteins}
}
\]

\[
\longrightarrow
\]

\[
\boxed{
\text{cells}
}
\]

\[
\longrightarrow
\]

\[
\boxed{
\text{tissues}
}
\]

\[
\longrightarrow
\]

\[
\boxed{
\text{organisms}
}
\]

\[
\longrightarrow
\]

\[
\boxed{
\text{populations}.
}
\]

A multiscale model attempts to connect processes operating on different
levels.

%%%%%%%%%%%%%%%%%%%%%%%%%%%%%%%%%%%%%%%%%%%%%%%%%%%%%%%%%%%%
\subsection{Multiple Biological Time Scales}
\label{subsec:biological-time-scales}
%%%%%%%%%%%%%%%%%%%%%%%%%%%%%%%%%%%%%%%%%%%%%%%%%%%%%%%%%%%%

Some biological variables evolve rapidly while others evolve slowly.

Suppose

\[
\boxed{
\varepsilon\frac{dx}{dt}
=
f(x,y),
}
\]

\[
\boxed{
\frac{dy}{dt}
=
g(x,y),
}
\]

with

\[
0<\varepsilon\ll1.
\]

Then \(x\) may evolve on a much faster time scale than \(y\).

Singular perturbation methods can exploit this separation.

%%%%%%%%%%%%%%%%%%%%%%%%%%%%%%%%%%%%%%%%%%%%%%%%%%%%%%%%%%%%
\subsection{Parameter Estimation in Biology}
\label{subsec:biological-parameter-estimation}
%%%%%%%%%%%%%%%%%%%%%%%%%%%%%%%%%%%%%%%%%%%%%%%%%%%%%%%%%%%%

Biological parameters may include:

\[
\boxed{
\text{growth rates},
}
\]

\[
\boxed{
\text{transmission rates},
}
\]

\[
\boxed{
\text{death rates},
}
\]

\[
\boxed{
\text{reaction constants},
}
\]

and

\[
\boxed{
\text{migration rates}.
}
\]

They are often estimated by fitting models to experimental or observational
data.

%%%%%%%%%%%%%%%%%%%%%%%%%%%%%%%%%%%%%%%%%%%%%%%%%%%%%%%%%%%%
\subsection{Biological Identifiability}
\label{subsec:biological-identifiability}
%%%%%%%%%%%%%%%%%%%%%%%%%%%%%%%%%%%%%%%%%%%%%%%%%%%%%%%%%%%%

A biological model may contain more parameters than available data can
determine.

For example, if an observed trajectory depends mainly on

\[
\frac{\beta}{\gamma},
\]

then \(\beta\) and \(\gamma\) may be difficult to estimate separately from
limited data.

Identifiability analysis should accompany parameter fitting.

%%%%%%%%%%%%%%%%%%%%%%%%%%%%%%%%%%%%%%%%%%%%%%%%%%%%%%%%%%%%
\subsection{Sensitivity in Biological Models}
\label{subsec:biological-sensitivity}
%%%%%%%%%%%%%%%%%%%%%%%%%%%%%%%%%%%%%%%%%%%%%%%%%%%%%%%%%%%%

Sensitivity analysis identifies parameters having strong influence on
outputs.

For model output \(Y\),

\[
\boxed{
S_i
=
\frac{\partial Y}{\partial\theta_i}
}
\]

measures local sensitivity to parameter \(\theta_i\).

This can guide:

\[
\boxed{
\text{experiments},
}
\]

\[
\boxed{
\text{data collection},
}
\]

and

\[
\boxed{
\text{intervention design}.
}
\]

%%%%%%%%%%%%%%%%%%%%%%%%%%%%%%%%%%%%%%%%%%%%%%%%%%%%%%%%%%%%
\subsection{Uncertainty in Biological Data}
\label{subsec:biological-data-uncertainty}
%%%%%%%%%%%%%%%%%%%%%%%%%%%%%%%%%%%%%%%%%%%%%%%%%%%%%%%%%%%%

Biological data may contain:

\[
\boxed{
\text{measurement noise},
}
\]

\[
\boxed{
\text{individual variability},
}
\]

\[
\boxed{
\text{sampling bias},
}
\]

\[
\boxed{
\text{missing observations},
}
\]

or

\[
\boxed{
\text{reporting delays}.
}
\]

Model conclusions should reflect these uncertainties.

%%%%%%%%%%%%%%%%%%%%%%%%%%%%%%%%%%%%%%%%%%%%%%%%%%%%%%%%%%%%
\subsection{Deterministic Versus Stochastic Biology}
\label{subsec:deterministic-stochastic-biology}
%%%%%%%%%%%%%%%%%%%%%%%%%%%%%%%%%%%%%%%%%%%%%%%%%%%%%%%%%%%%

Deterministic models often approximate average behavior in large
populations.

Stochastic models become especially important when:

\[
\boxed{
\text{populations are small},
}
\]

\[
\boxed{
\text{events are rare},
}
\]

or

\[
\boxed{
\text{random extinction matters}.
}
\]

The appropriate model depends on scale and purpose.

%%%%%%%%%%%%%%%%%%%%%%%%%%%%%%%%%%%%%%%%%%%%%%%%%%%%%%%%%%%%
\subsection{Continuous Versus Discrete Populations}
\label{subsec:continuous-discrete-populations}
%%%%%%%%%%%%%%%%%%%%%%%%%%%%%%%%%%%%%%%%%%%%%%%%%%%%%%%%%%%%

Differential-equation models treat population variables as continuous.

This may be reasonable when populations are large.

For very small populations,

\[
\boxed{
P=2.4
}
\]

individuals has no literal biological meaning.

Discrete stochastic models may then be more appropriate.

%%%%%%%%%%%%%%%%%%%%%%%%%%%%%%%%%%%%%%%%%%%%%%%%%%%%%%%%%%%%
\subsection{Model Validation in Biology}
\label{subsec:biological-model-validation}
%%%%%%%%%%%%%%%%%%%%%%%%%%%%%%%%%%%%%%%%%%%%%%%%%%%%%%%%%%%%

Biological models should be compared with independent observations whenever
possible.

Validation may involve:

\[
\boxed{
\text{time-series prediction},
}
\]

\[
\boxed{
\text{spatial patterns},
}
\]

\[
\boxed{
\text{dose-response curves},
}
\]

or

\[
\boxed{
\text{experimental interventions}.
}
\]

Agreement with one dataset does not guarantee general biological validity.

%%%%%%%%%%%%%%%%%%%%%%%%%%%%%%%%%%%%%%%%%%%%%%%%%%%%%%%%%%%%
\subsection{Biological Mechanism Versus Statistical Fit}
\label{subsec:biology-mechanism-fit}
%%%%%%%%%%%%%%%%%%%%%%%%%%%%%%%%%%%%%%%%%%%%%%%%%%%%%%%%%%%%

A model may fit biological data for the wrong mechanistic reason.

Different mechanisms can generate similar observable trajectories.

Therefore one should distinguish:

\[
\boxed{
\text{parameter fit}
}
\]

from

\[
\boxed{
\text{mechanistic validation}.
}
\]

Experiments designed to distinguish competing mechanisms can be especially
valuable.

%%%%%%%%%%%%%%%%%%%%%%%%%%%%%%%%%%%%%%%%%%%%%%%%%%%%%%%%%%%%
\subsection{Model Complexity in Biology}
\label{subsec:biological-model-complexity}
%%%%%%%%%%%%%%%%%%%%%%%%%%%%%%%%%%%%%%%%%%%%%%%%%%%%%%%%%%%%

Biological systems are extraordinarily complex.

A model should not include every known biological detail.

Useful modeling asks:

\[
\boxed{
\text{Which mechanisms are essential for the current question?}
}
\]

Adding unnecessary mechanisms can reduce interpretability and create
unidentifiable parameters.

%%%%%%%%%%%%%%%%%%%%%%%%%%%%%%%%%%%%%%%%%%%%%%%%%%%%%%%%%%%%
\subsection{Biological Modeling Workflow}
\label{subsec:biological-modeling-workflow}
%%%%%%%%%%%%%%%%%%%%%%%%%%%%%%%%%%%%%%%%%%%%%%%%%%%%%%%%%%%%

A useful workflow is:

\begin{enumerate}
    \item define the biological question;
    \item determine the biological scale;
    \item identify relevant populations, species, or molecular components;
    \item choose state variables;
    \item define biological parameters;
    \item identify conserved quantities where possible;
    \item specify interactions and transition mechanisms;
    \item determine whether time is discrete or continuous;
    \item determine whether randomness is important;
    \item determine whether spatial structure is important;
    \item formulate the governing equations;
    \item verify units;
    \item determine equilibria;
    \item analyze stability;
    \item identify biological thresholds;
    \item estimate parameters from data;
    \item test structural and practical identifiability;
    \item perform sensitivity analysis;
    \item quantify uncertainty;
    \item validate against independent biological observations;
    \item compare alternative mechanisms;
    \item communicate assumptions and limitations;
    \item revise the model when biological evidence requires it.
\end{enumerate}

%%%%%%%%%%%%%%%%%%%%%%%%%%%%%%%%%%%%%%%%%%%%%%%%%%%%%%%%%%%%
\subsection{Common Errors}
\label{subsec:mathematical-biology-common-errors}
%%%%%%%%%%%%%%%%%%%%%%%%%%%%%%%%%%%%%%%%%%%%%%%%%%%%%%%%%%%%

\subsubsection{Assuming Exponential Growth Can Continue Indefinitely}

Exponential growth ignores resource limitation.

It is often useful only over limited ranges.

%%%%%%%%%%%%%%%%%%%%%%%%%%%%%%%%%%%%%%%%%%%%%%%%%%%%%%%%%%%%
\subsubsection{Interpreting Carrying Capacity as an Immutable Constant}

The parameter \(K\) summarizes environmental constraints.

Real carrying capacity may change with:

\[
\boxed{
\text{climate},
}
\]

\[
\boxed{
\text{resources},
}
\]

or

\[
\boxed{
\text{habitat}.
}
\]

%%%%%%%%%%%%%%%%%%%%%%%%%%%%%%%%%%%%%%%%%%%%%%%%%%%%%%%%%%%%
\subsubsection{Assuming Lotka--Volterra Is a Complete Ecological Model}

The classical predator-prey equations omit:

\[
\boxed{
\text{predator saturation},
}
\]

\[
\boxed{
\text{prey carrying capacity},
}
\]

\[
\boxed{
\text{age structure},
}
\]

and many other ecological effects.

They are a foundational model, not a universal description.

%%%%%%%%%%%%%%%%%%%%%%%%%%%%%%%%%%%%%%%%%%%%%%%%%%%%%%%%%%%%
\subsubsection{Confusing \(R_0\) With the Current Number of Infections}

\(R_0\) is a reproduction number defined under a specific epidemiological
model and reference population state.

It is not the number of currently infectious individuals.

%%%%%%%%%%%%%%%%%%%%%%%%%%%%%%%%%%%%%%%%%%%%%%%%%%%%%%%%%%%%
\subsubsection{Treating \(R_0\) as Universal Across Models}

The formula

\[
R_0=\frac{\beta}{\gamma}
\]

belongs to the basic SIR framework.

More complex models can have different \(R_0\) formulas.

%%%%%%%%%%%%%%%%%%%%%%%%%%%%%%%%%%%%%%%%%%%%%%%%%%%%%%%%%%%%
\subsubsection{Confusing \(R_0\) and \(R_{\mathrm{eff}}\)}

\(R_0\) describes transmission potential in the reference susceptible
population.

The effective reproduction number changes with current immunity,
susceptibility, interventions, and other model factors.

%%%%%%%%%%%%%%%%%%%%%%%%%%%%%%%%%%%%%%%%%%%%%%%%%%%%%%%%%%%%
\subsubsection{Using \(\beta SI\) and \(\beta SI/N\) Interchangeably}

These incidence functions represent different transmission assumptions.

Their \(\beta\) parameters also have different units.

%%%%%%%%%%%%%%%%%%%%%%%%%%%%%%%%%%%%%%%%%%%%%%%%%%%%%%%%%%%%
\subsubsection{Forgetting Population Conservation}

For a closed SIR model,

\[
S+I+R=N.
\]

If numerical simulations violate this substantially, the implementation or
solver should be checked.

%%%%%%%%%%%%%%%%%%%%%%%%%%%%%%%%%%%%%%%%%%%%%%%%%%%%%%%%%%%%
\subsubsection{Treating Compartment Boundaries as Literal Biological Categories}

Compartments are modeling abstractions.

Real individuals may have heterogeneous disease stages and transition times.

%%%%%%%%%%%%%%%%%%%%%%%%%%%%%%%%%%%%%%%%%%%%%%%%%%%%%%%%%%%%
\subsubsection{Assuming All Waiting Times Are Exponential}

Ordinary compartment ODEs with constant transition rates imply exponential
waiting-time distributions.

Biological durations may have different distributions.

Additional stages, delays, or non-Markovian models may be required.

%%%%%%%%%%%%%%%%%%%%%%%%%%%%%%%%%%%%%%%%%%%%%%%%%%%%%%%%%%%%
\subsubsection{Ignoring Age or Stage Structure}

Birth, death, infection, and reproduction rates can vary strongly by age or
stage.

Aggregate models may fail when these differences are important.

%%%%%%%%%%%%%%%%%%%%%%%%%%%%%%%%%%%%%%%%%%%%%%%%%%%%%%%%%%%%
\subsubsection{Ignoring Small-Population Randomness}

A deterministic model may predict smooth persistence while a stochastic
population can become extinct.

%%%%%%%%%%%%%%%%%%%%%%%%%%%%%%%%%%%%%%%%%%%%%%%%%%%%%%%%%%%%
\subsubsection{Treating Model Parameters as Directly Observed Facts}

Parameters such as

\[
\beta
\]

or

\[
K
\]

may summarize several biological mechanisms.

Their interpretation depends on the model.

%%%%%%%%%%%%%%%%%%%%%%%%%%%%%%%%%%%%%%%%%%%%%%%%%%%%%%%%%%%%
\subsubsection{Ignoring Identifiability}

A numerical optimizer can return parameter values even when many different
parameter sets produce nearly identical biological predictions.

%%%%%%%%%%%%%%%%%%%%%%%%%%%%%%%%%%%%%%%%%%%%%%%%%%%%%%%%%%%%
\subsubsection{Confusing Correlation With Mechanism}

A fitted association between biological variables does not establish a
causal biological pathway.

%%%%%%%%%%%%%%%%%%%%%%%%%%%%%%%%%%%%%%%%%%%%%%%%%%%%%%%%%%%%
\subsubsection{Ignoring Units}

Transmission, growth, transfer, and reaction parameters have units determined
by the equations in which they appear.

Unit consistency should always be checked.

%%%%%%%%%%%%%%%%%%%%%%%%%%%%%%%%%%%%%%%%%%%%%%%%%%%%%%%%%%%%
\subsubsection{Ignoring Spatial Structure}

A well-mixed model may fail when populations are geographically separated or
interactions depend strongly on location.

%%%%%%%%%%%%%%%%%%%%%%%%%%%%%%%%%%%%%%%%%%%%%%%%%%%%%%%%%%%%
\subsubsection{Ignoring Network Structure}

Two populations with the same average number of contacts can have different
epidemic behavior if their network structures differ.

%%%%%%%%%%%%%%%%%%%%%%%%%%%%%%%%%%%%%%%%%%%%%%%%%%%%%%%%%%%%
\subsubsection{Assuming More Biological Detail Always Improves the Model}

Additional detail can increase:

\[
\boxed{
\text{uncertainty},
}
\]

\[
\boxed{
\text{computational cost},
}
\]

and

\[
\boxed{
\text{nonidentifiability}.
}
\]

Model complexity should serve the biological question.

%%%%%%%%%%%%%%%%%%%%%%%%%%%%%%%%%%%%%%%%%%%%%%%%%%%%%%%%%%%%
\subsubsection{Overinterpreting Numerical Precision}

Biological parameters may be uncertain even when a numerical solver returns
many decimal places.

Reported precision should reflect data and model uncertainty.

%%%%%%%%%%%%%%%%%%%%%%%%%%%%%%%%%%%%%%%%%%%%%%%%%%%%%%%%%%%%
\subsection{Applications}
\label{subsec:mathematical-biology-applications}
%%%%%%%%%%%%%%%%%%%%%%%%%%%%%%%%%%%%%%%%%%%%%%%%%%%%%%%%%%%%

\begin{applicationbox}

\textbf{Population Ecology.}
Growth, harvesting, competition, predator-prey interactions, migration, and
extinction can be studied using dynamical and stochastic models.

\medskip

\textbf{Epidemiology.}
SIR, SEIR, SEIRD, network, spatial, and stochastic models help analyze
disease transmission and interventions.

\medskip

\textbf{Pharmacology.}
Compartment models describe drug absorption, distribution, metabolism, and
elimination.

\medskip

\textbf{Biochemistry.}
Mass-action kinetics and enzyme models describe chemical reactions inside
living systems.

\medskip

\textbf{Genetics.}
Allele-frequency dynamics, mutation, selection, and inheritance can be
modeled using probability and dynamical systems.

\medskip

\textbf{Evolution.}
Replicator equations and evolutionary game theory model competition among
strategies and phenotypes.

\medskip

\textbf{Neuroscience.}
Nonlinear dynamical systems describe membrane potentials, signaling,
oscillations, and excitability.

\medskip

\textbf{Cancer Modeling.}
Growth laws, treatment models, evolutionary dynamics, and spatial models can
describe tumor progression.

\medskip

\textbf{Systems Biology.}
Networks of genes, proteins, metabolites, and signaling pathways are modeled
with differential equations, stochastic processes, and graphs.

\end{applicationbox}

%%%%%%%%%%%%%%%%%%%%%%%%%%%%%%%%%%%%%%%%%%%%%%%%%%%%%%%%%%%%
\subsection{Exercises}
\label{subsec:mathematical-biology-exercises}
%%%%%%%%%%%%%%%%%%%%%%%%%%%%%%%%%%%%%%%%%%%%%%%%%%%%%%%%%%%%

\begin{exercise}[1]
Define mathematical biology.
\end{exercise}

\begin{exercise}[1]
Define exponential population growth.
\end{exercise}

\begin{exercise}[1]
Define carrying capacity.
\end{exercise}

\begin{exercise}[1]
Write the logistic growth equation.
\end{exercise}

\begin{exercise}[1]
Define a predator-prey model.
\end{exercise}

\begin{exercise}[1]
Define an epidemic compartment.
\end{exercise}

\begin{exercise}[1]
Write the basic SIR model.
\end{exercise}

\begin{exercise}[1]
Define \(R_0\) conceptually.
\end{exercise}

\begin{exercise}[1]
Define a Leslie matrix.
\end{exercise}

\begin{exercise}[1]
Define mass-action kinetics.
\end{exercise}

\begin{exercise}[2]
Solve

\[
\frac{dP}{dt}=0.03P,
\qquad
P(0)=500.
\]
\end{exercise}

\begin{exercise}[2]
Find the doubling time for continuous growth rate

\[
r=0.08.
\]
\end{exercise}

\begin{exercise}[2]
Find the equilibria of

\[
\frac{dP}{dt}
=
rP
\left(
1-\frac{P}{K}
\right).
\]
\end{exercise}

\begin{exercise}[2]
Determine the sign of logistic growth for

\[
P<K
\]

and

\[
P>K.
\]
\end{exercise}

\begin{exercise}[3]
Show that logistic growth is maximized at

\[
P=\frac{K}{2}.
\]
\end{exercise}

\begin{exercise}[3]
Find the maximum logistic growth rate.
\end{exercise}

\begin{exercise}[3]
For the harvested logistic equation

\[
\frac{dP}{dt}
=
rP
\left(
1-\frac{P}{K}
\right)
-
H,
\]

determine when positive equilibria can exist.
\end{exercise}

\begin{exercise}[2]
Write the Lotka--Volterra predator-prey system.
\end{exercise}

\begin{exercise}[3]
Find the positive equilibrium of

\[
\dot{x}
=
\alpha x-\beta xy,
\qquad
\dot{y}
=
\delta xy-\gamma y.
\]
\end{exercise}

\begin{exercise}[3]
Interpret each Lotka--Volterra parameter biologically.
\end{exercise}

\begin{exercise}[3]
Explain the feedback cycle producing predator-prey oscillations.
\end{exercise}

\begin{exercise}[3]
Write a two-species competition model.
\end{exercise}

\begin{exercise}[3]
Explain what the coefficients

\[
\alpha_{12},
\qquad
\alpha_{21}
\]

represent.
\end{exercise}

\begin{exercise}[2]
Verify that the SIR model conserves total population.
\end{exercise}

\begin{exercise}[2]
For

\[
\beta=0.24,
\qquad
\gamma=0.08,
\]

compute \(R_0\).
\end{exercise}

\begin{exercise}[3]
Determine whether infection initially grows when

\[
R_0=2.5
\]

and

\[
\frac{S}{N}=0.3.
\]
\end{exercise}

\begin{exercise}[3]
Derive the condition for the peak of \(I(t)\) in the SIR model.
\end{exercise}

\begin{exercise}[3]
Derive the idealized herd-immunity threshold for the basic SIR model.
\end{exercise}

\begin{exercise}[2]
Write the SEIR equations.
\end{exercise}

\begin{exercise}[2]
Explain the meaning of \(\sigma\) in an SEIR model.
\end{exercise}

\begin{exercise}[3]
Extend the SEIR system to include disease-related mortality.
\end{exercise}

\begin{exercise}[3]
Add waning immunity to an SIR model.
\end{exercise}

\begin{exercise}[3]
Add continuous vaccination to an SIR model.
\end{exercise}

\begin{exercise}[3]
Compare frequency-dependent incidence

\[
\beta\frac{SI}{N}
\]

with density-dependent incidence

\[
\beta SI.
\]
\end{exercise}

\begin{exercise}[3]
Determine the units of \(\beta\) in both incidence formulations.
\end{exercise}

\begin{exercise}[3]
Explain why births can permit endemic infection.
\end{exercise}

\begin{exercise}[2]
Construct a three-age-class Leslie matrix.
\end{exercise}

\begin{exercise}[2]
Compute one projection step for a supplied Leslie matrix.
\end{exercise}

\begin{exercise}[3]
Explain the biological meaning of the dominant eigenvalue.
\end{exercise}

\begin{exercise}[3]
Explain the stable age distribution.
\end{exercise}

\begin{exercise}[3]
Construct a stage-structured projection matrix for seed, juvenile, and adult
plants.
\end{exercise}

\begin{exercise}[3]
Explain the continuous age-structured PDE

\[
\frac{\partial n}{\partial t}
+
\frac{\partial n}{\partial a}
=
-\mu(a)n.
\]
\end{exercise}

\begin{exercise}[2]
Solve

\[
\frac{dA}{dt}=-kA,
\qquad
A(0)=A_0.
\]
\end{exercise}

\begin{exercise}[2]
Find the drug half-life when

\[
k=0.15
\]

per hour.
\end{exercise}

\begin{exercise}[3]
Construct a two-compartment pharmacokinetic model.
\end{exercise}

\begin{exercise}[2]
For reaction

\[
A+B\to C,
\]

write the mass-action reaction rate.
\end{exercise}

\begin{exercise}[3]
Write differential equations for the concentrations in

\[
A+B\to C.
\]
\end{exercise}

\begin{exercise}[2]
Write the Michaelis--Menten equation.
\end{exercise}

\begin{exercise}[2]
Show that when

\[
S=K_M,
\]

the reaction rate is

\[
\frac{V_{\max}}{2}.
\]
\end{exercise}

\begin{exercise}[3]
Find the low-substrate approximation to the Michaelis--Menten equation.
\end{exercise}

\begin{exercise}[3]
Find the high-substrate approximation.
\end{exercise}

\begin{exercise}[2]
Write a Hill activation function.
\end{exercise}

\begin{exercise}[3]
Explain how the Hill coefficient affects response steepness.
\end{exercise}

\begin{exercise}[3]
Construct a basic mRNA-protein production model.
\end{exercise}

\begin{exercise}[3]
Add negative autoregulation using a Hill repression term.
\end{exercise}

\begin{exercise}[3]
Explain how positive feedback can generate bistability.
\end{exercise}

\begin{exercise}[2]
Write a generic reaction-diffusion system for two biological species.
\end{exercise}

\begin{exercise}[3]
Explain why diffusion can affect stability.
\end{exercise}

\begin{exercise}[3]
Construct a logistic reaction-diffusion population model.
\end{exercise}

\begin{exercise}[3]
Construct a network migration model for three habitat patches.
\end{exercise}

\begin{exercise}[3]
Explain the difference between a metapopulation model and one well-mixed
population model.
\end{exercise}

\begin{exercise}[2]
Define a birth-death process.
\end{exercise}

\begin{exercise}[3]
For

\[
\lambda_n=\lambda n,
\qquad
\mu_n=\mu n,
\]

interpret the rates biologically.
\end{exercise}

\begin{exercise}[3]
Explain why stochastic extinction can occur even when deterministic growth
is positive.
\end{exercise}

\begin{exercise}[3]
State the basic extinction criterion for a Galton--Watson process
conceptually.
\end{exercise}

\begin{exercise}[2]
Write the replicator equation.
\end{exercise}

\begin{exercise}[3]
Explain the role of mean fitness in the replicator equation.
\end{exercise}

\begin{exercise}[3]
Show that a strategy with fitness above the mean increases in frequency.
\end{exercise}

\begin{exercise}[3]
Construct a two-strategy evolutionary game from a \(2\times2\) payoff
matrix.
\end{exercise}

\begin{exercise}[2]
State the Hardy--Weinberg genotype frequencies.
\end{exercise}

\begin{exercise}[2]
For

\[
p=0.6,
\qquad
q=0.4,
\]

compute the genotype frequencies.
\end{exercise}

\begin{exercise}[3]
Construct a mutation-only allele-frequency recursion.
\end{exercise}

\begin{exercise}[3]
Explain how natural selection alters allele frequencies.
\end{exercise}

\begin{exercise}[2]
Write a logistic tumor-growth model.
\end{exercise}

\begin{exercise}[2]
Write a Gompertz growth model.
\end{exercise}

\begin{exercise}[3]
Add treatment-induced mortality to a tumor-growth model.
\end{exercise}

\begin{exercise}[3]
Construct a simple pathogen-immune interaction model.
\end{exercise}

\begin{exercise}[3]
Explain how contact networks modify epidemic models.
\end{exercise}

\begin{exercise}[3]
Give an example of a biological network other than an epidemic network.
\end{exercise}

\begin{exercise}[3]
Explain biological multiscale modeling.
\end{exercise}

\begin{exercise}[3]
Give an example of fast and slow variables in a biological model.
\end{exercise}

\begin{exercise}[3]
Explain structural identifiability in a biological context.
\end{exercise}

\begin{exercise}[3]
Explain practical identifiability.
\end{exercise}

\begin{exercise}[3]
Explain how sensitivity analysis can guide experimental design.
\end{exercise}

\begin{exercise}[3]
List several sources of uncertainty in biological data.
\end{exercise}

\begin{exercise}[3]
Explain why biological model validation should use independent evidence when
possible.
\end{exercise}

\begin{exercise}[3]
Explain why a model can fit data well while representing the wrong
mechanism.
\end{exercise}

\begin{exercise}[3]
Explain when a stochastic biological model may be preferable to a
deterministic one.
\end{exercise}

\begin{exercise}[4]
Analyze the stability of the logistic equilibrium \(P=K\).
\end{exercise}

\begin{exercise}[4]
Analyze the Jacobian of the Lotka--Volterra system at its coexistence
equilibrium.
\end{exercise}

\begin{exercise}[4]
Study a predator-prey model with logistic prey growth.
\end{exercise}

\begin{exercise}[4]
Study a saturating functional response in a predator-prey model.
\end{exercise}

\begin{exercise}[4]
Analyze coexistence in a two-species competition system.
\end{exercise}

\begin{exercise}[4]
Derive \(R_0\) using the next-generation matrix method for a simple
multi-compartment epidemic system.
\end{exercise}

\begin{exercise}[4]
Construct an SEIRS model with births and natural deaths.
\end{exercise}

\begin{exercise}[4]
Analyze the disease-free equilibrium of an epidemic model.
\end{exercise}

\begin{exercise}[4]
Construct a multi-region epidemic model with travel.
\end{exercise}

\begin{exercise}[4]
Study an age-structured epidemic model.
\end{exercise}

\begin{exercise}[4]
Analyze the dominant eigenvalue of a Leslie matrix.
\end{exercise}

\begin{exercise}[4]
Develop a continuous age-structured population model.
\end{exercise}

\begin{exercise}[4]
Derive the Michaelis--Menten approximation from the enzyme-reaction system.
\end{exercise}

\begin{exercise}[4]
Analyze a gene-regulatory feedback system.
\end{exercise}

\begin{exercise}[4]
Study bistability generated by positive feedback.
\end{exercise}

\begin{exercise}[4]
Study a reaction-diffusion pattern-formation model.
\end{exercise}

\begin{exercise}[4]
Construct a stochastic birth-death simulation.
\end{exercise}

\begin{exercise}[4]
Compare stochastic and deterministic logistic growth.
\end{exercise}

\begin{exercise}[4]
Analyze a two-strategy replicator system.
\end{exercise}

\begin{exercise}[4]
Study mutation-selection balance.
\end{exercise}

\begin{exercise}[4]
Compare logistic and Gompertz tumor-growth models.
\end{exercise}

\begin{exercise}[4]
Perform parameter estimation for a biological ODE model.
\end{exercise}

\begin{exercise}[4]
Perform sensitivity analysis on an epidemic or ecological model.
\end{exercise}

\begin{exercise}[4]
Develop a spatial metapopulation model on a graph.
\end{exercise}

%%%%%%%%%%%%%%%%%%%%%%%%%%%%%%%%%%%%%%%%%%%%%%%%%%%%%%%%%%%%
\subsection{Conceptual Summary}
\label{subsec:mathematical-biology-summary}
%%%%%%%%%%%%%%%%%%%%%%%%%%%%%%%%%%%%%%%%%%%%%%%%%%%%%%%%%%%%

Mathematical biology uses mathematics to represent biological mechanisms
across scales.

Population growth begins with models such as

\[
\boxed{
\frac{dP}{dt}=rP
}
\]

and

\[
\boxed{
\frac{dP}{dt}
=
rP
\left(
1-\frac{P}{K}
\right).
}
\]

Ecological interactions lead to coupled nonlinear systems such as

\[
\boxed{
\dot{x}
=
\alpha x-\beta xy,
}
\]

\[
\boxed{
\dot{y}
=
\delta xy-\gamma y.
}
\]

Epidemic models divide populations into compartments such as

\[
\boxed{
S,
E,
I,
R,
D.
}
\]

For the basic SIR model,

\[
\boxed{
R_0
=
\frac{\beta}{\gamma}
}
\]

provides an epidemic threshold under the model assumptions.

Population structure can be represented through matrix models such as

\[
\boxed{
\mathbf{n}_{t+1}
=
L\mathbf{n}_t.
}
\]

Biochemical dynamics use reaction laws such as

\[
\boxed{
v=k[A][B]
}
\]

and enzyme models such as

\[
\boxed{
v
=
\frac{V_{\max}S}{K_M+S}.
}
\]

Spatial biology introduces diffusion:

\[
\boxed{
\frac{\partial u}{\partial t}
=
D\nabla^2u
+
f(u).
}
\]

Evolutionary dynamics can be described by

\[
\boxed{
\dot{x}_i
=
x_i
(f_i-\bar f).
}
\]

Across all these models, the essential modeling questions remain:

\[
\boxed{
\text{Which mechanisms matter?}
}
\]

\[
\boxed{
\text{Which parameters can be estimated?}
}
\]

\[
\boxed{
\text{How sensitive are the conclusions?}
}
\]

and

\[
\boxed{
\text{Does the model agree with biological evidence?}
}
\]

Mathematical biology therefore combines

\[
\boxed{
\text{differential equations}
+
\text{probability}
+
\text{statistics}
+
\text{networks}
+
\text{computation}
}
\]

to study living systems quantitatively.

%%%%%%%%%%%%%%%%%%%%%%%%%%%%%%%%%%%%%%%%%%%%%%%%%%%%%%%%%%%%
\subsection{Section Connections}
\label{subsec:mathematical-biology-connections}
%%%%%%%%%%%%%%%%%%%%%%%%%%%%%%%%%%%%%%%%%%%%%%%%%%%%%%%%%%%%

\chapterconnection{
Mathematical Biology applies the general modeling framework of
Section~\ref{sec:mathematical-modeling} to populations, ecosystems,
epidemics, biochemical reactions, genetics, evolution, physiology, and
cellular systems. Differential equations describe biological dynamics,
probability captures demographic and molecular randomness, linear algebra
supports structured populations, networks represent interaction patterns,
statistics supports parameter estimation, and numerical methods make
complex models computationally accessible. The next section, Mathematical
Physics, develops mathematical models built from physical laws, including
mechanics, fields, waves, variational principles, conservation laws,
relativity, and quantum systems.
}

%%%%%%%%%%%%%%%%%%%%%%%%%%%%%%%%%%%%%%%%%%%%%%%%%%%%%%%%%%%%
% SECTION SOURCE TRACKING
%%%%%%%%%%%%%%%%%%%%%%%%%%%%%%%%%%%%%%%%%%%%%%%%%%%%%%%%%%%%

%------------------------------------------------------------
% 13.2 Mathematical Biology
%------------------------------------------------------------
%
% Source basis:
%   - Standard mathematical biology.
%   - Standard population dynamics, epidemiology, ecology,
%     pharmacokinetics, biochemical kinetics, evolutionary dynamics,
%     genetics, and biological-network modeling.
%
% Used for:
%   - population growth;
%   - exponential growth;
%   - logistic growth;
%   - carrying capacity;
%   - harvesting;
%   - predator-prey systems;
%   - Lotka--Volterra equations;
%   - competition;
%   - mutualism;
%   - epidemic compartment models;
%   - SIR;
%   - SEIR;
%   - SEIRD;
%   - reproduction numbers;
%   - epidemic thresholds;
%   - effective reproduction numbers;
%   - vaccination;
%   - waning immunity;
%   - frequency- and density-dependent transmission;
%   - vital dynamics;
%   - endemic equilibria;
%   - age-structured populations;
%   - Leslie matrices;
%   - stage structure;
%   - continuous age structure;
%   - pharmacokinetic models;
%   - drug half-life;
%   - biochemical reaction networks;
%   - mass-action kinetics;
%   - Michaelis--Menten kinetics;
%   - Hill functions;
%   - gene regulation;
%   - neural models;
%   - excitable systems;
%   - reaction-diffusion systems;
%   - spatial ecology;
%   - migration;
%   - metapopulations;
%   - stochastic birth-death models;
%   - branching processes;
%   - extinction;
%   - evolutionary dynamics;
%   - replicator equations;
%   - evolutionary game theory;
%   - population genetics;
%   - Hardy--Weinberg equilibrium;
%   - mutation and selection;
%   - cancer-growth models;
%   - immune-response models;
%   - biological networks;
%   - contact networks;
%   - multiscale biological modeling;
%   - multiple time scales;
%   - parameter estimation;
%   - identifiability;
%   - sensitivity analysis;
%   - uncertainty;
%   - model validation;
%   - applications and exercises.
%
% Notes:
%   - Mathematical Physics is developed next in Section 13.3.
%   - Continuum Mechanics follows in Section 13.4.
%   - Fluid Dynamics follows in Section 13.5.
%   - Control Theory follows in Section 13.6.
%   - Network Science follows in Section 13.7.
%
%%%%%%%%%%%%%%%%%%%%%%%%%%%%%%%%%%%%%%%%%%%%%%%%%%%%%%%%%%%%